% 原书第 269 页（PDF 第 282 页）
\BookPage{269}{282}
\BookHeading{康熙小臣养心殿总监造赵昌生平小考}{康熙小臣养心殿总监造赵昌生平小考}{陈国栋*}
「赵昌」这个名字在《熙朝定案》、《正教奉褒》等中国天主教文献中频繁出现，而早在雍正九年（1731），法国耶稣会传教士冯秉正（Joseph-François-Marie-Anne de Moyriac de Mailla）就在报导赵昌受洗的一封信函中，介绍过这个人。\BookRef{1}因此在天主教的著作中，赵昌也时而被述及。\BookRef{2}但在已往，由于文献不容易到手，因此聚焦于这个小人物的研究不多，直到1996年《康熙朝满文朱批奏折全译》\BookRef{3}出版以后，相关研究中才较常提到赵昌，例如冯明珠在论述中梵关系史的一篇文章中就大量征引《康熙朝满文朱批奏折全译》，并且十分注意赵昌在为康熙皇帝传话给西洋教士一事上所扮演的重要角色。\BookRef{4}更加值得注意的是在过去三、四年间，竟然出现了两篇论文以赵昌为唯一的研究对象，并且揭露了一些过去罕为人知的事实。此一课题的发展，在某种程度上，反映了最近这一、二十年来，研究资源的大量公开与取得的相对便利，有助于特定议题的开发与研究。

\vfill
\begin{BookNotes}
\BookNote{*}{作者任职于中央研究院历史语言研究所。}
\BookNote{1.}{耿升译，《耶稣会士中国书简集》，第4册（河南郑州：大象出版社，2005），页66-74，〈耶稣会士冯秉正神父致同一耶稣会的法国传教区会长赫苍璧神父的信〉，发信的地点在北京。该信函所押的日期为1731年10月10日，耿升的译本误植为1751年10月10日。该信的内容主要在报导赵昌的下场与天主教士设法安排他受洗的经过。参考矢泽利彦编译，《イエズス會士中國書簡集（第6冊—信仰篇）》（东京：平凡社，1974），页245-259所收本篇译文。}
\BookNote{2.}{例如河北省献县天主堂1931年原刊的萧若瑟，《天主教传行中国考》，便对赵昌受洗一事大加着墨，见辅仁大学编，《中国天主教史籍汇编》（台北县新庄市：辅仁大学出版社，2003），页200。方豪，〈清代旗人之信奉天主教与遭禁〉一文第四小节，标题即为「侍卫赵昌之信教」，见《故宫文献》，第4卷第4期（1973年9月），页6-7。方豪在此一节文章中，主要为整理《熙朝定案》与《正教奉褒》二书中与赵昌有关的事迹，但他没有引用到冯秉正的信函。}
\BookNote{3.}{中国第一历史档案馆译编，《康熙朝满文朱批奏折全译》（北京：中国社会科学出版社，1996）。}
\BookNote{4.}{冯明珠，〈坚持与容忍—档案中所见康熙皇帝对中梵关系生变的因应〉，收在天主教辅仁大学历史系编辑，《中梵外交关系史国际学术研讨会论文集》（台北县新庄市：辅大历史系，2002），页145-182；页169-181的附表有系统地整理了康熙朝中国与教廷交涉相关的奏折档案，极有参考价值。戴逸在两篇论述雍正皇帝的文章中也简略地论及赵昌，参考戴逸，〈历史上的雍正〉、〈雍正继位的历史疑谜〉，《戴逸自选集》（北京：中国人民大学出版社，2007），页60, 69-70。}
\end{BookNotes}

% 原书第 270 页（PDF 第 283 页）
\BookPage{270}{283}
较早正式出版的一篇是天津南开大学的博士生陈青松的作品〈赵昌家世及其与传教士的往来一兼述其在康雍时期的际遇〉，刊登在韩国庆北大学亚细亚研究所发行的期刊《亚洲研究》第六期（2009年8月），很容易从网路读取。

另一篇的作者是旅居葡萄牙的金国平（JinGuoping）先生。2008年11月，他为在里斯本及澳门两地分别举行的会议写了一篇文章。会议的名称是：「红墙内外一徐日升神父（1645-1708），康熙皇帝与耶稣会士在华传教工作」；金国平先生的论文题目是〈情同手足—赵昌与徐日升(Amicíssimos- Tomás Pereira e Zhao Chang）。金国平在会议中发表的论文系用葡萄牙文写作，会后所准备的英文稿目前正在由牛津大学出版社进行编校，不日将可问世；此外，金国平也将原稿改写为中文，增补后于2010年在北京中国人民大学所举办的清史会议上（与吴志良联名）发表，题目为〈西方史料所记载的赵昌〉，近期也将出版。作者承金先生好意，惠赐校对中稿本一份，得享先睹为快之乐，谨向金先生致谢。

陈青松的论文主要建立在中文以及译汉文献（含满文译汉）的开发上；金国平的中文稿当然已经参考了陈青松的作品，但他更进一步的贡献则是征引了大量的欧洲语文资料。经由掌握前所未见的丰富史料，两位作者在史实重建上都作出了优异的成绩，并且在若干地方也提出清楚有见地的解释。

个人在资料方面只能依循两位作者的方向继续扩充，稍稍拾遗补阙而已；而基于论述的需要，也无法避免重复引用两位作者作品中已经征引过的文献，乃至于转引他们所开发出来的史料。不同之处是，本文主要尝试的方向放在将相关资料作进一步的考订与解读，设法厘清某些关键性的问题。经此处理之后，希望能对赵昌的生平作出更为完整的叙述，借以了解他与康熙皇帝的君臣关系。不过，赵昌毕竟只是个「小臣」，最后以被雍正皇帝下狱抄家，了结一生，因此能在史料中留下的纪录往往也只是一鳞半爪，要完整重建他的故事颇难，其间更不免有所猜测。望祈读者包涵，将本文作为进一步讨论之基础，暂时容忍其不够坚实的部份论证。

\vfill
\begin{BookNotes}
\BookNote{5.}{金国平的文章正待刊出，页码不确定，引用其文章时，将仅提其大名，不另行加注。}
\end{BookNotes}

% 原书第 271 页（PDF 第 284 页）
\BookPage{271}{284}
\BookSubheading{一、赵昌的基本资料}

赵昌的一生都在宫廷工作。他是一位内务府的「包衣」（满文booi，家人或家仆的意思）。就身体情况而言，包衣是一般的正常人。不过，少数人倒误以为赵昌是名太监，或许是因为看到他经常随侍在康熙皇帝左右的缘故。

陈青松认为赵昌的旗籍不详，但从他在被捕之初，被枷号于东直门示众这点，依据八旗在北京城派住的方位推断，「正白旗在东直门内」，认为他很可能是正白旗人，怀疑枷号于东直门，也许「有让他在自家门口丢人现眼的目的」。8其实，依「八旗方位」而言，上三旗的正黄旗居西直门内，而正白旗与镶黄旗分别居于东直门两侧。若依陈青松的逻辑，赵昌也有可能属于镶黄旗。

有一个时间极为相近的个案可供参考，即雍正元年二月二十五日，刑部尚书宗室佛格上奏请旨，事关是否可以释放一名被枷号的人犯。该奏折提到，雍正元年正月初九日，挑选鹰场拜唐阿（按、即鹰匠）时，被引见给皇帝的候选者当中，有一位镶黄旗的鹰匠索柱，系官员之子，竟于康熙皇帝大丧期间不着缨帽而身穿着绿袄，颇为不敬，因此奉上谕将他拿捕，送至东直门枷号。但因为眼前皇帝又有恩诏释放戴枷之人，于是佛格就为索柱奏请是否也能被释放？得旨鞭一百下后释放。既有恩诏释放戴枷之人了，为什么佛格等人还要上折奏请呢？这是因为索柱是「奉特旨枷号之人，不敢援恩诏释放」，所以另行请旨。这位被枷示东直门的鹰匠正好是镶黄旗人。那么，我们是否可以援引这个事例推断赵昌也是镶黄旗人呢？恐怕并不这么简单。

就制度而言，北京城在清初设立步军统领一职，康熙十三年兼提督九门事务，

\vfill
\begin{BookNotes}
\BookNote{6.}{请参考陈国栋，〈清代内务府包衣三旗人员的分类及其旗下组织〉，《食货月刊》，第12卷第9期(1982年12月)，页5-23。}
\BookNote{7.}{例如杨启樵，〈刘心武先生的《揭秘红楼梦》质疑〉，《红楼梦学刊》，2007年第4辑，页304就直接说「康熙亲信太监赵昌」；参考冯明珠，〈坚持与容忍—档案中所见康熙皇帝对中梵关系生变的因应〉，页149-150。}
\BookNote{8.}{陈青松，〈赵昌家世及其与传教士的往来—兼述其在康雍时期的际遇〉，页4，注3。}
\BookNote{9.}{佛格（等），〈刑部奏为枷号人犯请旨释放折〉，见河内良弘，〈满汉合璧雍正朝奏折译注〉，《京都大学文学部研究纪要》，第31辑（1992)，页65-66；参考庄吉发《雍正朝满汉合璧奏折校注（第1辑）》（台北：文史哲出版社，1984），页21-22。}
\end{BookNotes}

% 原书第 272 页（PDF 第 285 页）
\BookPage{272}{285}
\noindent 「审理八旗（巡捕）三营拏获违禁犯法奸匪逃盗一应案件」。10

\begin{BookQuote}
雍正二年议准：崇文、宣武、朝阳、阜成、东直、西直、安定、德胜八门，每门各设监狱。凡旗人获罪，刑部议定枷示，或步军统领衙门奏明永远枷示人犯，皆发门监羇禁。
\end{BookQuote}

并且明定「崇文门监镶白旗人，宣武门监镶红旗人，朝阳门监正蓝旗人，阜成门监镶蓝旗人，东直门监镶黄旗人，西直门监正黄旗人，安定门监正白旗人，德胜门监正红旗人。」11并不以本旗所在地界的门监为收禁、枷示该旗犯人为原则。虽然依此规定，镶黄旗人会被安排在东直门枷示，与索柱的个案正好相符，不过，在索柱案发生当时，尙未议准收禁枷示各旗旗人的地点，而且索柱是「奉特旨枷号之人」，也不是通例。赵昌之被捕及枷示发生于索柱之前，他更是「奉特旨枷号之人」，因此就算是被枷示于东直门，也无法据此推论说他是正白旗、镶黄旗或者正黄旗中那一旗的包衣人。12

探求赵昌的旗分，另有一种办法是查阅乾隆九年（1744）编定的《八旗满洲氏族通谱》。不过，在该书的汉文本当中，我们找不到「赵昌」这个名字。换一个方式，如果我们能确定赵昌的满洲名字的话，还可以利用满文本来检查。

康熙四十四年（1705）教宗使臣安提阿宗主教多罗（Carlo TommasoMaillard de Tournon,Patriarch of Antioch)到达北京，激化了所谓的「中国礼仪」之争。13康熙皇帝派出艾若瑟(Antonio FrancescoGiuseppeProvana)神父前往罗马交涉，但长久没有消息。于是皇帝在康熙五十五年（1716）下令以满文、汉文及拉丁文三体文字，以朱色印刷一张可以称为公开信的「红票」（RedManifesto），循各种管道送往欧洲，以迫使教廷作出回应。赵昌列名在上面，因此我们可以比对在不同语

\vfill
\begin{BookNotes}
\BookNote{10.}{参考胡祥雨，〈清前期京师初级审判制度之变迁〉，《历史档案》，2007年第2期，页42-54,60。}
\BookNote{11.}{以上见昆冈等修，《钦定大清会典事例》（台北:台湾中文书局，1963），卷1158，页1（p. 18635)。九门当中不设门监者为正阳门。}
\BookNote{12.}{「康熙十一年奏准：议处内务府官员有与吏、兵二部颁定律例相符者，即照依吏部、兵部律例定议。见昆冈等修，《钦定大清会典事例》，卷1212，页1（p.19148)。依此，索柱虽然是旗分人，而赵昌是包衣人，在雍正即位之初的法律待遇应该相似。}
\BookNote{13.}{参考安双成，〈礼仪之争与康熙皇帝（上）、（下）〉，《历史档案》，2001年第1期，页74-80；2001年第2期，页32-41。也请参考本文第五节。}
\end{BookNotes}

% 原书第 273 页（PDF 第 286 页）
\BookPage{273}{286}
\noindent 文中他的名字的写法。我们发现在拉丁文版的人名“Tchaotcham”，显然即是汉语「赵昌」两字的音译。这与现代汉语的拼法十分近似。14不过，「红票」上的满文，以罗马字拼出来则是“Joocang”。15

我们利用“Joocang”这个满文拼音到满文本《八旗满洲氏族通谱》里寻找，可以找到四位同名者，对照他们在汉文本应该出现的地方，这四个人都被写作「兆昌」。16旗人人名汉译无定字，显然「兆昌」应该就是「赵昌」的另一种写法。

四名「兆昌」在商务印书馆翻印的文渊阁本汉文本《八旗满洲氏族通谱》17，分别出现在15/8a,15/8b、26/15a以及75/4a。其中，15/8a、15/8b的两位「兆昌」皆属「各地方伊尔根觉罗氏」，前者为镶蓝旗旗分人，后者为正红旗旗分人。出现在26/15a的那位「兆昌」属讷殷地方富察氏，正白旗旗分人。内务府成员基本上是上三旗包衣人，以上三位「兆昌」皆为旗分人，因此不可能是本文所讨论的赵昌。《八旗满洲氏族通谱》最后几卷为「附载满洲旗分内之尼堪姓氏」，也就是汉姓人的简谱。出现在75/4a的兆昌姓王，很可能就是我们的赵昌。

\begin{BookQuote}
《八旗满洲氏族通谱》相关的文字云：
\end{BookQuote}

\begin{BookQuote}
王盘，正黄旗包衣管领下人，世居沈阳地方。天聪时来归。其子王受贵原任蓝翎侍卫。曾孙兆昌，原任库掌。元孙常住，现任笔帖式。
\end{BookQuote}

这里所载的兆昌「原任库掌」。「库长」在满文中称为「乌林达」（ulin-i da），官品不定，但常为六品。我们虽然没有本文主人翁赵昌担任库长的资料，而赵昌的官衔变来变去，但在绝大部份时间里，他都任职于养心殿造办处，为六品官。因此被记录为「原任库长」，还不算背离事实。

\vfill
\begin{BookNotes}
\BookNote{14.}{威妥马式拼法拼为“ChaoCh'ang”；汉语拼音拼为“Zhao Chang”，。}
\BookNote{15.}{「红票」上全部的满文文字，德国学者Walravens在其作品内已将之完全用罗马字拼出。参考Hartmut Walravens, "Zu Zwei Katholischen Katechismen in Mandjurischer Sprache," Monumenta Serica, vol. 31 (1977 for 1974-76), pp. 522-523.}
\BookNote{16.}{参考 Giovanni Stary, A Dictionary of Manchu Names: a name-index to the Manchu version of the "Complete genealogies of the Manchu clans and families ofthe eight banners"; jakūn gūsai Manjusai mukūn hala be uheri ejehe bithe; Baqi manzhou shizu tongpu;八旗满洲氏族通谱(Wiesbaden :Harrassowitz,2000)。}
\BookNote{17.}{台北:台湾商务印书馆，1983。}
\end{BookNotes}

% 原书第 274 页（PDF 第 287 页）
\BookPage{274}{287}
\begin{BookQuote}
《八旗满洲氏族通谱》的记载说明了他的曾祖父是王盘，原籍沈阳，天聪
\end{BookQuote}

(1627-1636)年间入旗，从龙入关。王盘之子王受贵，应该也就是赵昌的祖父，原任蓝翎侍卫。《八旗满洲氏族通谱》完全没有提到赵昌父亲这一辈的任何名字。不过，据《张诚日记》所记，1688年时，赵昌的父亲仍然健在，并且也出内务府的差。据此推断，其在内务府当差所达到的位置恐怕只是苏拉（闲散)之流而已。因为未能获得任何职名，所以《八旗满洲氏族通谱》也就跳过他，不记录他的名字。

不过，陈青松与金国平都注意到耶稣会神父冯秉正在那封报导赵昌受洗的信函中提到了赵昌的父亲。陈青松的相关文字如下：

据冯（秉正）信称，「赵老爷（赵昌）是康熙皇帝之父顺治皇帝侍从中的一位

\begin{BookQuote}
一品大员的儿子。」18
\end{BookQuote}

金国平从同一封信函的原文本加以引用，对中文文句略作修饰，他的文稿说：

\begin{BookQuote}
冯秉正神父言：「赵老爷是康熙皇帝之父顺治帝随从中的一名一等大员的儿子(Tchao-laoye est fils d’un des grands du premier ordre qui étoient à la suite de Chun-tchi, père du feu empereur Cang-hi.)。」19其父余况不明。
\end{BookQuote}

法文的“undes grandsdupremier ordre”的确可以译成「一品大员」。但是从文献中所见赵昌的生平行止，他实在不可能有个一品大官的父亲，否则也不会连在《八旗满洲氏族通谱》留下名字的机会都没有。事实上，在《耶稣会士中国书简集》的译本中，冯秉正还这样子形容赵昌这个人。他说：

\begin{BookQuote}
赵老爷是康熙皇帝之父顺治皇帝侍从中的一位一品大员的儿子。当他尚很年幼时，便是宫中业绩最佳者之一，以其天生美男子的仪表，思维敏捷，举止文雅和行为审愼而成为所有青年王爷中的佼佼者。20
\end{BookQuote}

\vfill
\begin{BookNotes}
\BookNote{18.}{陈青松，〈赵昌家世及其与传教士的往来—兼述其在康雍时期的际遇〉，页4。陈青松使用的版本即是耿升所译的〈耶稣会士冯秉正神父致同一耶稣会的法国传教区会长赫苍璧神父的信〉，页66。}
\BookNote{19.}{原注出处为Carta doP. de Moyriac de Mailla ao P. Hervieu(冯秉正神父致赫苍璧神父函)，Pequim,}
\BookNote{10.}{10.1731, Lettres Édifiantes et Curieuses Écrites des Missions Étrangères, ed. de Lyon, 1819, tomo. XII, p.88.两种译法都忽略了“feu”这个字一「最近亡故的」之意。}
\BookNote{20.}{〈耶稣会士冯秉正神父致同一耶稣会的法国传教区会长赫苍璧神父的信〉，页66。}
\end{BookNotes}

% 原书第 275 页（PDF 第 288 页）
\BookPage{275}{288}
\noindent 竟然用「王爷」来称呼赵昌，可以说非常夸大不严谨。赵昌虽然经常随侍在皇帝左右，但几乎一直都只是个六品小官而已。在身分上，他是管领下包衣人，也不会是具有高尙地位的人。

陈青松的文章还提了有关赵昌之父的一点点讯息。他说：

\begin{BookQuote}
根据张诚1688年的旅行日记，至少在当时赵昌之父仍然在世，当时因为张诚与徐日升的帐篷尚未运到，他俩还去赵父的驻处，并且受到热情接待。21
\end{BookQuote}

此外，关于赵昌之父也就没有进一步线索可寻了。

赵昌的祖父王受贵为蓝翎侍卫。「蓝翎侍卫」为皇宫侍卫最低的一个等级，较高的等级分别叫作一等侍卫、二等侍卫和三等侍卫。蓝翎侍卫位阶不高，很少高于六品。22

到得康熙六十一年十一月，新即位的雍正皇帝逮捕赵昌，将他枷号示众后下狱。逮捕他后，于抄家时发现他的祖父坟墓的守墓人共有十一口。陈青松和金国平都依据乾隆《钦定大清会典则例》的规定，推定赵昌的家世还不错。个人对相关文献的理解略有不同。先来看一下有关赵昌祖父守墓人的那段文字。

汇报抄家结果的文件当中，有一件雍正元年，月日不详的折片云：「查得：守赵昌祖父之坟共有男女大小共十一口。」对此雍正皇帝以一种宽大的口吻写下「将此赏之」四个字的朱批。23

陈青松和金国平都提到本条，并且引用了《钦定大清会典则例》的资料来说明。《钦定大清会典则例》为乾隆十二年（1747）奉敕撰成的典籍，其卷一百六十为有关内务府会计司职掌的政令与事例。当中特别有「振恤」（按、即赈恤）一条，其文字云：

\vfill
\begin{BookNotes}
\BookNote{21.}{陈青松，〈赵昌家世及其与传教士的往来一兼述其在康雍时期的际遇〉，页4。}
\BookNote{22.}{参考《清史稿》卷118，〈职官五：内务府〉；秦国经，〈清代宫廷的警卫制度〉，《故宫博物院院刊》，1990年第4期，页65。}
\BookNote{23.}{内务府奏报追还赵昌欠银办法片〉，收在中国第一历史档案馆译编，《雍正朝满文殊批奏折全译》（合肥：黄山书社，1998），上册，页591。}
\end{BookNotes}

% 原书第 276 页（PDF 第 289 页）
\BookPage{276}{289}
\begin{BookQuote}
康熙九年奏准（按、准）：各佐领、内管领下绝嗣之人，有奴仆田地者给守墓二人、祭田三十亩，仍归本佐领、内管领下。有奴仆、无田地者，给予守墓二人外，余人交部收管。或有亲生之女，亦给祭田如前数。内务府三旗绝嗣产业，有亲女、亲姊妹者，除守墓奴仆、祭田外，其余人口、田产交该佐领、内管领分给亲女二分、姊妹一分。但有姊妹者，除给守墓奴仆、祭田外，系（按、系）同佐领、内管领者，并令承受人口、家产；系（按、系）别佐领、内管领者，止令承受奴仆、器皿，其田产房屋交与该旗。24
\end{BookQuote}

在将赵昌下狱，查抄赵昌家产，并将其家口进行处置之时，有可能赵昌的父、祖已被视同为内务府佐领或内管领下的「绝嗣之人」。赵昌既被抄家，奴仆、财产应行充公，应该算是一无所有了。不过，在获罪之前，他本来是有奴仆、田地之人家。若从后者认定，则依康熙九年的规定，可以为其祖父之坟墓留下两名守墓人，并且保留三十亩地作为祭田。当时，依据内务府的调查，守赵昌祖父之坟者共有男女大小口十一人，雍正皇帝慷慨地让他们全部继续留下来守墓。至于祭田部份，因为资料不足，无法知晓。赵昌在长久服务与经营之后，原本家道殷实，是京师里公认的有钱人。他的祖父之坟有十一名守墓人也不见得多，因为连绝户都被认为要分配给他们两名守墓人了，更何况赵昌祖父去世应该已有数十年之久，这十一个人很可能是原守墓人结婚娶妇以及生养后代才达到的数字，最初可能也只有少数几位。陈松青与金国平两位都从会计司的这条规定推论说「赵昌先祖的地位不低」，其实不甚相干，因为只要家有奴仆或田产的内务府家庭绝嗣时，都可以留下守墓人两名，原户主的身分并未纳入考虑。从条文称作「振恤」来看，这是一种慈善之举的考量。

赵昌被查抄时有妻妾、儿子，也有儿媳及孙辈。署理内务府总管事务允禄、办理内务府总管事务领侍卫内大臣曼武、内务府总管噶达浑、办理内务府总管事务佐领伊都立负责抄家，他们在雍正元年正月初六日〈内务府奏查赵昌家产事折〉中报告雍正皇帝说：

\vfill
\begin{BookNotes}
\BookNote{24.}{《钦定大清会典则例（六）》（台北市:台湾商务印书馆，1983，「景印文渊阁四库全书」，第625册），页202-203。}
\end{BookNotes}

% 原书第 277 页（PDF 第 290 页）
\BookPage{277}{290}
\begin{BookQuote}
内务府谨奏，为遵旨查赵昌家产事。康熙六十一年十二月初二日25奉上谕：除赵昌之妻及儿媳外，将闲散女人内现有子者查后另奏。将无子之闲散女人赏与包衣佐领及管领内无妻贫穷之人。钦此。钦遵。经派员外郎瓦尔达、原看守王道华26查核，除将赵昌之闲散女人内，三名无子之闲散女人，赏给管领内贫穷之人外，赵昌夫妻、子九人、儿媳四人、孙三人、孙女二人，算有子之闲散女人三人，共二十三口。其有子之三名闲散女人之年龄、姓名，已缮入另一折片内。27
\end{BookQuote}

另有一件雍正元年失月日折片:〈内务府奏报赵昌有子之闲散女人数片〉，显然就是上引正月初六日奏折内所说的折片，报告赵昌的「有子之闲散女人数」情形如下：

\begin{BookQuote}
赵昌有子之闲散女人三：鄂尔杰，五十四岁。其子显宁，二十六岁。秦杰，四十九岁。其子三保，十七岁；勒山，十五岁；仁寿，十三岁。苏伦朱，四十岁。其子山庆，十九岁。28
\end{BookQuote}

除了妻妾子女等亲属之外，赵昌还拥有相当数量的奴婢和动产、不动产。雍正元年正月初六日，〈内务府奏查赵昌家产事折〉29有巨细靡遗的记录。重要的项目包括家内奴婢一百一十一口，在庄园上的男女及其他各处的奴婢共一百零八口、房屋五百多间、田地五十六顷七十六亩一分。在动产方面，各种珠宝、衣饰、武器、马、骡、驼、牛等项目和数量都相当可观。特别值得注意的是，他还拥有：

\vfill
\begin{BookNotes}
\BookNote{25.}{康熙皇帝于六十一年十一月十三日驾崩，距此不足二十日，而赵昌已然下狱抄家多日，可知他在十一月间就已经被捕。}
\BookNote{26.}{此人当即是王道化，在康熙后期担任养心殿总监造，是赵昌的旧同事。参见《康熙朝满文朱批奏折全译》，页405，〈广储司银皮库奏报装饰养心殿所用物件数目折〉（康熙四十五年正月二十二日）。参考方豪，《中国天主教史人物传》（北京：宗教文化出版社，2007），页417-423及韩琦，〈康熙时代的历算活动：基于档案资料的新研究〉，收在张先清编，《史料与视界》（上海：上海人民出版社，2007），页40-60。}
\BookNote{27.}{《雍正朝满文殊批奏折全译》，上册，页2。}
\BookNote{28.}{《雍正朝满文殊批奏折全译》，上册，页591。}
\BookNote{29.}{《雍正朝满文殊批奏折全译》，上册，页2-4。}
\end{BookNotes}

% 原书第 278 页（PDF 第 291 页）
\BookPage{278}{291}
\begin{BookQuote}
大小钟表六；…；玻璃镜大小十三；各种玻璃小物件一百九十二；各种西洋物品一百六十八种；大小千里眼六…30
\end{BookQuote}

这些进口的钟表、玻璃制品、各种西洋物品…等，不是一般人有机会拥有的物件，赵昌之所以会拥有，系与他长期与西洋人往来有关，我们留在第五节再进一步讨论。

就一位贴身为皇帝工作的人物而言，赵昌被查抄的财产当中，其实不乏他当差任职的行头。整体而言，也不见得过分或逾制，因此他也没有因为拥有清单上的任何物品而被安上新的罪名。他算是富有，事实上康熙皇帝也说过赵昌是「名富」（有名的富人），但这一切都是在他服务康熙皇帝五、六十年后的情况。31至于他为何富裕，多少是因为职务的关系而获得赀财的缘故吧。当然，他也从事经营。康熙四十九年三月，他就曾经从内库以年利百分之十的代价借出了五千两银子。虽然不知用途，但可能用于投资生利。32金国平引用了多位葡萄牙神父的报导说赵昌还将资金寄到澳门投资。

不过，赵昌充其量也只是有钱而已。作为管领下人的他，在内务府内的地位不高。许多有名的同时代的内务府包衣（如曹寅、赫世亨…等）都担任过内务府的旗务官佐领（牛录章京）或内管领（包衣大），赵昌却没有担任过佐领、内管领等类似的职务。

\vfill
\begin{BookNotes}
\BookNote{30.}{内务府总管们建议「将钟表交自鸣钟修造处，将玻璃器皿等物交烧玻璃处。」见同一奏折。}
\BookNote{31.}{康熙皇帝之所以知道赵昌是「名富」，多少是当时还实行满洲传统的「派巴延」制度有关。该制度为在当时的八旗制度下，掌握近期发财或长期家赀富有的成员资料，在政府或皇室有特殊需要的时候，派给这些人去负责担当。为图顺利交差，这些人不免贴钱办事，而皇帝或政府、八旗倒节省了一笔开销。为此作法，同时也建立了一些掌握旗人财产状况的办法，所以康熙皇帝很容易知道赵昌是位「名富」。关于「派巴延」制度，请参考陈国栋，〈清代中叶以后重要税差专由内务府包衣担任的几点解释〉，收入许倬云、毛汉光、刘翠溶主编，《第二届中国社会经济史研讨会论文集》(台北：汉学研究资料及服务中心，1983)，页173-204，特别是页182，注30。}
\BookNote{32.}{《雍正朝满文朱批奏折全译》，上册，页591。}
\end{BookNotes}

% 原书第 279 页（PDF 第 292 页）
\BookPage{279}{292}
\BookSubheading{二、赵昌的其他名字：Pursai与Chuliama}

赵昌与西洋人，特别是传教士的接触可以说是他最为人知晓的事情，因此在《熙朝定案》、《正教奉褒》或者《康熙朝满文朱批奏折全译》…等文献中都能看到相关的实例。也由于这样频繁的接触，他常被西洋人加以记录、常出现在传教士的文献当中。奇怪的是，除了赵昌、赵大人…等这类正常的称呼之外，据欧、美学者的研究，他在传教士文献中还被称作Pursai或Chuliama。

先说Pursai吧。孟德卫(DavidMungello)在以下二处就说赵昌是Pursai:

第一处是康熙二十八年（1689）第二次南巡，阳历2月，皇帝在济南遇到传教士柯若瑟（JosephdeOsca），召见他。柯若瑟说他不懂数学与天文，中文也不够好，所以虽然来到中国已经两年了，都还没去过北京。皇帝叫他返回教堂，说他的侍卫赵昌（原注:Pursai)可能会在稍后去拜访他。33

第二处发生的时间是在几天之后的2月28日，地点则由济南换到杭州。康熙南巡，在杭州遇到传教士殷铎泽(ProsperoIntorcetta)，召见他。当殷铎泽与皇帝的谈话触及杨光先的当儿，有一位名叫赵昌（原注:即Pursai）的内务府满洲官员扬声说:虽然他经常造访天主堂，他对他所见的一切一概不明了。两天之后，赵昌带着另外五名侍卫造访了殷铎泽的教堂。34

在这两个地方，孟德卫都直接把赵昌与Pursai认定为是同一个人的不同名字。《正教奉褒》在相关的日期下也记载了皇帝在济南会见柯若瑟与在杭州会见殷铎泽的事情。提到赵昌时只称呼为「赵侍卫」，完全没有提到Pursai。至于赵昌强调他去天主堂但完全不了解他所见到的一切这样的内容，《正教奉褒》中也完全没有提及。3值得注意的是赵昌在皇帝面前（御前）竟能扬声发语，显然是得到皇帝的默

\vfill
\begin{BookNotes}
\BookNote{33.}{David E. Mungello, The Spirit and the Flesh in Shandong, 1650-1785 (Lanham: Rowman \& Littlefield Publishers,2001),p.68.赵昌亦见p.71；本次南巡中，赵昌代替皇帝造访了许多教堂。}
\BookNote{34.}{David E. Mungello, The Forgotten Christians of Hangzhou(Honolulu: University of Hawaii Press, 1994) , P}
\BookNote{60.}{}
\BookNote{35.}{黄伯禄，《正教奉褒》，收在韩琦、吴旻校注，《熙朝崇正集、熙朝定案（外三种)》（北京:中华书局，2006)，页346-348。}
\end{BookNotes}

% 原书第 280 页（PDF 第 293 页）
\BookPage{280}{293}
\noindent 许。

中文（汉语）与Pursai接近的写法应是「布尔赛」或「布尔赛」。中文文献中，得确也有以这样名字出现的资料。

第一件为史语所内阁大库所藏明清档案第104464号，为一敕书，不完整，内容为「敕宗人府副理官布尔赛监督北新关事务」，时间为康熙三十三年六月初八日，满汉合璧，盖「敕命之宝」印两次。

第二件为雍正元年九月二十日「管理工部事务和硕廉亲王允裸」等人的奏折，所奏的事由为「臣部自康熙二十六年起至五十八年七月止题参案内，郎中马克察等四百零八人，共欠架木五十二万八千八百二十二根。」「再原任木仓监督员外郎布尔赛等十一人，共缺少桅木二百八十三根、杉木三千一百六十六根半。…」勒限分赔追结。「其员外郎布尔赛等十四人所欠桅杉木植，只有布尔赛缴回桅木二根，尚欠桅杉木三千四百九十九根半。伊等逾限二年有余，并未交还，殊属不合。」36

另外，康熙五十三年五月十二日，理藩院的奏折也提到一位布尔赛：

\begin{BookQuote}
理藩院谨奏，为请旨事：据住喀尔喀公颜楚布多尔济旗照看之员外郎布尔赛、员外郎扎西等呈称：去年闰五月37，经具奏将我等派往。我等来业已一年，请将我换回。等语。查得：先前臣等部奏称，往喀尔喀公颜楚布多尔济旗下派员照看时，各部衙门派出满洲章京一员，臣等部派出蒙古章京一员，遣往，满一年更换。等语，在案。今布尔赛、扎西等既已去一年，伏乞将伊等换回。.38
\end{BookQuote}

从「宗人府副理官监督北新关事务」、「木仓监督员外郎」与「（理藩院）员外郎」…等职称来推断，都与赵昌的其他资料不相符，以上三个文件所提到的布尔赛应该与赵昌不是同一个人。

\vfill
\begin{BookNotes}
\BookNote{36.}{〈和硕廉亲王允禩等奏参拖欠架木等案折〉，《雍正朝满文朱批奏折全译》，上册，页359-360。}
\BookNote{37.}{康熙五十二年闰五月，即西元1713年6月23日至7月21日。}
\BookNote{38.}{〈理藩院奏请更换派往喀尔喀公颜楚布多尔济旗之官员折〉，《康熙朝满文殊批奏折全译》，页943-944。}
\end{BookNotes}

% 原书第 281 页（PDF 第 294 页）
\BookPage{281}{294}
南开大学历史学院的白新良考察康熙朝进折人员与折件的下落，指出有一位「武英殿监督布尔赛」曾经单独或领衔进过奏折，但目前已经佚失。而据其说明，该项资讯的来源是《康熙朝满文朱批奏折全译》第1610件。39经查原书该处，并无此项讯息。正确的出处是《康熙朝汉文朱批奏折汇编》第1610件。该件为四川巡抚年羹尧的奏折，上奏时间是康熙五十三年七月初二日，折末提到「布尔赛等奏折一封内五件一并附进」。年羹尧奏折称呼布尔赛的职衔用的是「武英殿监视」五个字，与不是白新良所记的「武英殿监督」。差一个字，不过其意思没有太大的差别。年羹尧奏折的主旨在于报告四川省配合办理测绘地图4°的情况。41

康熙晚期命令西洋传教士分赴全国各地测绘地图是一件极为有名的事情。42正式开始的时间为康熙四十七年四月十六日（1708年6月4日）。皇帝命令法国传教士白晋(JoachimBouvet)、雷孝思（Jean-BaptisteRégis)，杜德美（Pierre Jartoux）和日耳曼神父费隐（Ehrenbert Xavier Fridelli）等人，从长城开始测量，而在次年一月完成一图。其后重组人员，测绘他处。整个测绘工作在康熙五十七年完成。43

测绘地图这件事与西洋人有关，而负责处理西洋人事务的就是武英殿和养心殿的官员。然而如后文所见，在整个测绘地图期间，赵昌都在养心殿总监造上任职。那么，他还可能是「武英殿监视布尔赛」吗？

康熙五十一年十二月十九日，武英殿与养心殿的监造们上了一个如下的奏折：

\begin{BookQuote}
和素、王道化、李国屏谨奏，为奏闻事。
\end{BookQuote}

\vfill
\begin{BookNotes}
\BookNote{39.}{白新良，〈康熙朝进折人员考略〉，《历史档案》，2005年第2期，页40；《康熙朝满文殊批奏折全译》，页685-686。}
\BookNote{40.}{即后来的《皇舆全览图》。}
\BookNote{41.}{中国第一历史档案馆编，《康熙朝汉文朱批奏折汇编》，第5册（北京市：档案出版社，1985），页674-676。}
\BookNote{42.}{参考方豪的三篇文章：〈康熙间测绘滇黔舆图考〉、〈康熙五十八年清廷派员测绘琉球地图之研究〉、〈康熙五十三年测绘台湾地图考〉，收在《方豪六十自订稿》（台北：著者自刊，1969)，上册，页515-604。}
\BookNote{43.}{秦国经，〈18世纪西洋人在测绘清朝舆图中的活动与贡献〉，《清史研究》，1997年1月，页37-}
\BookNote{44。}{}
\end{BookNotes}

% 原书第 282 页（PDF 第 295 页）
\BookPage{282}{295}
\begin{BookQuote}
本月十二日，兵部笔帖式玛索色将陕西省画图向导、护军参领觉罗英珠等带来之奏折一，带给奴才之文书一，送来后，因文书系带给赵昌、张常柱和奴才者，故奴才等于本月十三日，除将英珠等之奏折、文书带给赵昌、张常住44外，又趁具奏英珠奏折之便，一并咨文，欲参奏佛保，并押解前来惩治。十八日夜所到之报内，赵昌、张常住咨称：如何办理，尔等定夺。等语。将折及信一并带回。佛保系遣往画地图之人，虽艰苦，本应竭力跟随行走到底，嗣后落后，名声又恶，故请将伊带回惩治。已将向导、护军参领觉罗英珠等之奏折带出，而所带之文书内称，若总督、巡抚不奏，将此所带之折，亦请停具奏。等因。甚为不合。故将英珠等之奏折及带给奴才等之文书一并谨奏。请旨。【朱批：】尔等此奏甚是，依议行。45
\end{BookQuote}

会衔上奏的三人当中，和素、李国屏皆为武英殿监造，王道化则是养心殿监造。文件中提到的赵昌属养心殿，张常柱（张常住）则为武英殿监造。当时康熙皇帝前往遵化谒陵，赵昌与张常住显然是随扈出巡，与留在北京的其他监造保持联系。46赵昌虽然不在北京，但绝对不是去做测绘地图的工作。

几个月之后，有一名「监视武英殿布尔赛」被派出参加测绘地图的工作。《明清史料》收录了康熙五十二年（1713）的一个文件，内容如下：

\begin{BookQuote}
一件。钦奉上谕事。康熙五十二年闰五月十一日，准兵部车驾清吏司案呈，准养心殿咨称：康熙五十二年四月二十九日，监视养心殿李秉忠47，监视武英殿布尔赛等八人；西洋
\end{BookQuote}

\vfill
\begin{BookNotes}
\BookNote{44.}{张常住也可以写成张常柱或张长柱。不过，在此一文件中，同一个人名，却译写成不一样的汉字，还是有点奇怪。}
\BookNote{45.}{《康熙朝满文殊批奏折全译》，页835-836，〈武英殿总监造和素等奏请将佛保押来惩治折〉。}
\BookNote{46.}{参考王思治、冯尔康（编着)，《康熙事典》（北京：紫禁城出版社，2010)，页408。}
\BookNote{47.}{李秉忠的职务其实就是养心殿总监造，此处只是用字不同而已。史景迁认为在康熙后期，李秉忠与赵昌、赫世亨三人在康熙皇帝与天主教人士之间扮演关键的角色。参考JonathanD.Spence,"Claims and Counter-Claims: The Kangxi Emperor and the Europeans (1661-1722) , in D. E. Mungello ed., The Chinese Rites Controversy: Its History and Meaning (Nettetal: Steyler Verlag, 1994) , p.23.}
\end{BookNotes}

% 原书第 283 页（PDF 第 296 页）
\BookPage{283}{296}
人费隐、单爻占48、麦大成49、汤尚贤50往四川等九省去，列名绿头牌启奏。奉旨：布尔赛、李秉忠着派出。四川、贵州、云南省，着布尔赛、西洋人费隐、单爻占去。回来，从湖广省画江西、广东、广西省。着李秉忠、西洋人麦大成、汤尚贤去河南、江南二省画。去的官员若画完，就从彼处往浙江、福建画去。再，云南、贵州、广西瘴气甚不好，去的人员于四月至八月正遇瘴气，暂且往别处去画，等待瘴气散的时候再画去。钦此。

据此，布尔赛与西洋人费隐、单爻占（山遥瞻，GuillaumeFabreBonjour)三个人的第一波任务是先往四川、贵州和云南三省进行测绘。同一个文件又引录了两天后的另一道上谕。该上谕提到为了细节的安排，布尔赛

于初四日，将向导护军参领郎图、郎中绰尔代、五官正51苗受等十三人列名绿头牌启奏，奉旨：…往山西、陕西省去的向导护军参领英柱，郎中郎乌里、监副双德，着往四川、云南、贵州、湖广去。再，下旨与画舆图去的人等，如画完一省，就交与该抚派的当家人好生送来。钦此。四川、云南、贵州、湖广省，去的西洋人费隐、单爻占各骑马五匹，监视武英殿布尔赛骑马四匹，驮仪器马一匹，向导护军参领英柱、郎中郎乌里、监副双德等各骑马五匹。…伊等绘画通省府州县卫所舆图，若无驿马之处，仍照先年给与营马乘骑。等因。前来。相应行文各该省督抚提镇，转行所属作速预备壮营马，俟画图各官一到，如无驿马之处，即照勘牌内马数拨给营马乘骑，毋至迟误可也。等因。案呈到部，合咨前去。查照施行。等因。.52

布尔赛、费隐、山遥瞻一行先到四川进行地图测绘。第二年夏天便已告成。康

\vfill
\begin{BookNotes}
\BookNote{48.}{即山遥瞻(GuillaumeFabre Bonjour)。}
\BookNote{49.}{João Francisco Cardoso°}
\BookNote{50.}{PierreVincente de Tartre°}
\BookNote{51.}{钦天监的职官，正六品或从六品。}
\BookNote{52.}{《明清史料·丁编》（台北：中央研究院历史语言研究所，1997)，第8本，页768，〈西洋人费隐等奉旨往四川等九省画地理图散叶〉；本文件被改题为〈西洋人费隐等奉旨往四川等九省画地理图记注〉收入中国第一历史档案馆、澳门基金会、暨南大学古籍研究所合编，《明清时期澳门问题档案文献汇编》（北京市：人民出版社，1999），页103-105。参考吴伯娅，《康雍干三帝与西学东渐》（北京：宗教文化出版社，2002)，页395-397。本文的标点与以上这两件资料略有不同。}
\end{BookNotes}

% 原书第 284 页（PDF 第 297 页）
\BookPage{284}{297}
\noindent 熙五十三年七月初二日四川巡抚年羹尧奏报：

\begin{BookQuote}
于康熙五十二年六月十二日，布尔赛等由四川北界保宁府属之广元县画起，于康熙五十三年六月初十日至东川府画完川省全图。即将全图交与臣家人魏之辉收明。布尔赛等随由东川入滇画图去讫。53
\end{BookQuote}

武英殿监督布尔赛、西洋人费隐等人应该就此进入云南继续工作了。据秦国经引用第一历史档案馆所藏康熙五十四年六月二十四日，云南巡抚甘国璧的奏折，最后是由「西洋人费隐、雷孝思及武英殿监督常保」等人完成云南舆图的测绘，由甘国璧派遣家人送京呈览。54显然布尔赛没有做完云南部份的工作，而他因何没有入滇继续测绘，情况也不明。在方豪有关康熙末年云南地图测绘的研究论文中，也不见布尔赛的踪影。55至于他是否被召回北京还是另有他故，个人也一无所悉。然而，从以上的证据，我们至少可以确定有一武英殿监督（或称监视），名字写作布尔赛或布尔赛的人，从康熙五十二年六月二十二日到次年六月初十日间，确实就陪着西洋人在四川从事绘制地图的工作。

这一位布尔赛会不会是赵昌呢？应该不是。一则是不管写作「武英殿监视」还是「武英殿监督」，职衔上带的都是「武英殿」之名而不是赵昌职称所冠的「养心殿」。再则是，既然康熙五十二年下半年布尔赛确定在四川，则不可能出现在别处。可是，第一历史档案馆所藏的档案当中，却有一件康熙五十二年十一月二十八日的奏折如下：

\begin{BookQuote}
赵昌、和素、王道化、李国屏谨奏：奉旨：交李煦送来。钦此。二人中居凤阳之弗朗西亚国人郭仲川56本月二十六日到。问其年龄，答曰四十九岁；问其学问，答曰：天文、律吕皆不会，中国五经、四书粗看过。等语。令认字，粗识些。说汉话，还清楚。
\end{BookQuote}

\vfill
\begin{BookNotes}
\BookNote{53.}{中国第一历史档案馆编，《康熙朝汉文朱批奏折汇编》，第5册，页675。}
\BookNote{54.}{秦国经，〈18世纪西洋人在测绘清朝舆图中的活动与贡献〉，页39。}
\BookNote{55.}{方豪，〈康熙间测绘滇黔舆图考〉，页715-721。}
\BookNote{56.}{即郭中传(Jean-Alexis de Gollet)。}
\end{BookNotes}

% 原书第 285 页（PDF 第 298 页）
\BookPage{285}{298}
\begin{BookQuote}
再，住建昌府之西洋人马国宪57仍未到，为此恭奏以闻。58
\end{BookQuote}

换言之，康熙五十二年十一月赵昌人在北京，领衔与其他武英殿监造、养心殿监造一起上奏，报告有关郭中传、马国贤等西洋人的事情。他人不在四川现场，不可能是布尔赛。

至于内务府的这位布尔赛，他以武英殿监督或监视的名义出现，似乎就只发生在康熙五十二年和五十三年间。不过，他的确服务于内务府，而且也和处理西洋人的事务有关。如在教廷使节多罗到访期间，康熙皇帝原先派遣白晋（Joachim Bouvet）与沙国安（SabinoMariani）前往欧洲致送答礼，依据第一历史档案馆所藏的「内务府满文行文档」在康熙四十四年十一月的一件奏折，

\begin{BookQuote}
十九日，总管内务府咨行兵部，本月十七日，据员外郎赫世亨、张长柱告称，我等奏称，笔帖式布尔赛、西洋人白晋、沙国安抵达福建后，若海船尚未行，俟西洋人等登上海船，布尔赛即行返回；倘若海船业已起航，布尔赛或返回，或至广州之处，请旨。等因。奏入，奉旨：海船若已开行，则送至广东澳门。倘澳门有海船，俟登船即回；倘无返回海船，广州府有天主堂，可将西洋人等交付彼处巡抚安置于天主堂，布尔赛自行返回，西洋人搭明年返回海船前往。返回之消息，由巡抚等转告尔等。钦此。...59
\end{BookQuote}

Rosso在其书中叙及此事时，直接说1706年1月4日（康熙四十四年十一月二十日），白晋和沙国安由赵昌陪同，离开北京，乘驿前往广州。60Rosso自注说：

\vfill
\begin{BookNotes}
\BookNote{57.}{即马国贤(MatteoRipa)。}
\BookNote{58.}{《康熙朝满文殊批奏折全译》，页918，〈赵昌等奏闻弗朗西亚国人郭仲川情形折〉。}
\BookNote{59.}{《明清时期澳门问题档案文献汇编》，页68-69，康熙四十四年十一月十九日（1706年1月3日），<总管内务府为护送西洋人白晋等搭船回国并派拨驿马送往赏物事致兵部咨文〉。覃波使用了这个文件，但误作康熙三十二年白晋奉康熙皇帝之命返回欧洲招募人才之事。时间与事由都弄错了。见覃波，〈康熙皇帝与法国远臣〉，《紫禁城》，2005年第3期，页35。}
\BookNote{60.}{Antonio Sisto Rosso,Apostolic Legations to China of the Eighteenth Century (South Pasadena: P. D. and IonePerkins,1948)，p.161；参考陈垣编，《康熙与罗马使节关系文书影印本》(北平：国立北平故宫博物院，1932;台北:文海，1974)，页1-2及页7。Francis A.Rouleau(胡天龙），"Maillard de Tournon, Papal Legate at the Court of Peking: The First Imperial Audience (31 December 1705) ,"Archivum HistoricumSocietatis Iesu, no.31（1962)，p. 290也引用了多罗的资料，但完全没有提到陪同的官员，只说快马加鞭的结果，从北京到厦门只花十三天就到了。}
\end{BookNotes}

% 原书第 286 页（PDF 第 299 页）
\BookPage{286}{299}
\noindent 赵昌的满洲名字是Pursai和Purgama；而在1706年3月28日白晋致沙国安的函件中，此一官员的名字被写成是Parsama。他认为这些名字都是指同一个人，也就是赵昌。61

由于白晋、沙国安没赶上船，而两人又因为身分的问题引起纠纷，最后一行人被召回北京。康熙四十五年十二月初九日（1707年1月12日）有一件满文奏折报导了事件的后续：

\begin{BookQuote}
赫世亨谨奏：窃布尔赛携西洋人伯金、沙国安及其赏物于本月初八日晚到。令沙国安住徐日升之堂。初九日晨，布尔赛入内，交付所携赏物毕，已起行往迎皇上。为此谨奏以闻。62
\end{BookQuote}

除非布尔赛在前一年没等到船，在安置了白晋与沙国安后就先行北返，否则从四十四年底到四十五年底，前后超过一年期间，他是不在北京的。但是至少在康熙四十五年七月间赵昌随侍在皇帝左右，人在北京一带。63

柯兰霓(Claudiavon Collani)所着的白晋传记也只说赵昌提议由白晋陪同沙国安担任使节，并没有说赵昌一路陪伴他们南下。4要言之，从以上资料看来，陪同白晋与沙国安前往广州的人确实是布尔赛，但不可能是赵昌。

或许是源自于Rosso的说法，在1995年捷克刊行的、康熙末年来华的耶稣会士严嘉乐（KarelSlavíček）的书信集中，编辑者就下了一条注说：

\begin{BookQuote}
…赵昌，在那个时代关于传教工作的文件常用他的满文名字，如Pursai、Purgama、Pursama等，不尽统一。他属于满清皇族成员，在为中国礼仪之争中联系到教皇派铎罗主教和嘉乐主教来中国等的文稿中常提及赵昌的名字。…〔从
\end{BookQuote}

\vfill
\begin{BookNotes}
\BookNote{61.}{Antonio Sisto Rosso, Apostolic Legations to China of the Eighteenth Century, p. 161 note 32 and p. 433.}
\BookNote{62.}{《康熙朝满文殊批奏折全译》，页476，〈武英殿总监造赫世亨奏报西洋人迎驾日期折〉。}
\BookNote{63.}{见《康熙朝满文朱批奏折全译》，页434所收奏折及殊批。}
\BookNote{64.}{柯兰霓(ClaudiavonCollani)着、李岩译，《耶稣会士白晋的生平与著作》（郑州：大象出版社，2009），页52。}
\end{BookNotes}

% 原书第 287 页（PDF 第 300 页）
\BookPage{287}{300}
\begin{BookQuote}
略〕\BookRef{65}
\end{BookQuote}

由于中译者没有将整条注都加以翻译，因此我们无法明确知悉原注者的依据。

除了 Pursai、Purgama、Pursama（Parsama）等名字之外，金国平于「纪念徐日升神父国际研讨会」论文摘要提到赵昌也被称为 Chuliama。在其提供给笔者参考的论文修订稿中则没有再提 Chuliama 一名，当然也没有说明其来由。金国平之外，德国学者柯兰霓在一篇讨论严当（Charles Maigrot）的文章中，也明确地以赵昌为 Chuliama。严当因为反对中国天主教徒敬拜祖先及孔子，导致「礼仪之争」的激化。康熙皇帝派了赫世亨及赵昌两人去找严当，要他写下他认为儒家经典中与天主教矛盾的文句。这在叙述这两位康熙皇帝派去的人物时，柯兰霓是这样写的：

\begin{BookQuote}
Kangxi sent the two mandarins Zhao Chang (Chuliama 赵昌) and Henkama (or Hesken, Wang Daohua 王道化) to Maigrot with the imperial order to write down and deliver a compilation of sentences from the Confucian books which he thought to be in contradiction with Christianity.\BookRef{66}
\end{BookQuote}

因为柯兰霓没有交待赵昌为何就是 Chuliama，我们无法继续讨论。可以确定的是，他将 Henkama 当成王道化，绝对是错的。\BookRef{67}王道化与赫世亨为两个不同的人，有一段期间同时担任武英殿总监造。\BookRef{68}Henkama 也就是 Hesken，正是赫世亨。赫世亨与赵昌共同担任此一任务，也见于这两个人会衔的满文奏折。\BookRef{69}

赵昌与西洋传教士往来频繁，在西洋文献中留下纪录的机会自然很多。但光是他的名字的叫法，可能就不只一种。以上的考证只能说明 Pursai、Purgama、Pursama、Parsama 或 Chuliama 之类的称呼很难一概比定为是赵昌，乃至于根本

\vfill
\begin{BookNotes}
\BookNote{65.}{严嘉乐着，丛林、李梅译，《中国来信（1716-1735）》（郑州：大象出版社，2002），〈严嘉乐从南昌寄给布拉格尤利乌斯・兹维克尔的信（1723年10月14日）〉，页163，注释77。}
\BookNote{66.}{Claudia von Collani, "Charles Maigrot's Role in the Chinese Rites Controversy," in D. E. Mungello ed., The Chinese Rites Controversy: Its History and Meaning, p. 162.}
\BookNote{67.}{不只是柯兰霓，更早的 Rosso 就已经误以王道化为赫世亨了。见 Antonio Sisto Rosso, Apostolic Legations to China of the Eighteenth Century, p. 157 note 23.}
\BookNote{68.}{《康熙朝满文朱批奏折全译》等文献经常可以看到他们两人同时出现。}
\BookNote{69.}{例见《康熙朝满文朱批奏折全译》，页419-420，康熙四十五年五月二十七日，赫世亨与赵昌联名奏上的〈武英殿总监造赫世亨奏报西洋人情形折〉；另见同书多处，不具引。}
\end{BookNotes}

% 原书第 288 页（PDF 第 301 页）
\BookPage{288}{301}
\noindent 不是。由于在未来更多人使用西洋文献时，不免遇到这些称呼，研究者仍须仔细详查，才能避免不必要的误认。

\BookSubheading{三、哈哈珠子的问题}

金国平的文稿利用各种欧洲语文资料，说明赵昌从小便与康熙皇帝作件，重建了精彩的史实。最有意思的是他也引用了一条十分重要的资料。他说：

\begin{BookQuote}
徐日升神父在1688年的一信中称：「赵是皇帝最早的侍童（agatjusé或称pagem de companhia)。」70
\end{BookQuote}

徐日升(Tomás Pereira)为葡萄牙籍耶稣会士。葡萄牙文“pagemde companhia”指的就是「小跟班」的意思。它在徐日升信函中也用同义词“agatjusé”来表达，该词正是徐日升所属时代常用的满文叫法「哈哈珠子」（hahajuse），小跟班男孩（littlepageboy）的意思。金国平的文稿中，引用了数件传教士的文件，多处提到赵昌小时候就开始为康熙皇帝服务，甚至于如安文思神父（GabrieldeMagalhaes）在1669年时所描述的，赵昌是「皇帝的内寝侍童(menino da camara doRey)」。71至于冯秉正则在报导赵昌受洗的那封信中说「当他尚很年幼时，便是宫中业绩最佳者之一。他成了那些人们最早选出来的青年皇帝的伴读人之一。」72

关于清代早期的「哈哈珠子」，福格在其《听雨丛谈》，卷十二，〈哈哈珠子〉条有云:

\begin{BookQuote}
国初大臣子弟以童穉入侍禁近者，不必尽是旗人。康熙中，商丘文康相国宋权
\end{BookQuote}

\vfill
\begin{BookNotes}
\BookNote{70.}{原注:Carta de 20.9.1686doP TomásPereira ao P. Sebastião de Magalhães(徐日升神父致塞巴斯蒂昂·德.马加良埃斯神父函); cópia do P. Domingos MaurícioS.J. de uma carta mss. que pertenceu ao P. Francisco Rodrigues S.J., historiador da Companhia de Jesus, Arquivo da Cúria Provincial, Província Portuguesa da Companhia de Jesus, Pasta"Padre Tomás Pereira S.J."}
\BookNote{71.}{原注 : Carta de 25/29.12.1668 de Gabriel de Magalhães aos Padres e Superiores de Macau e Cantãoda Companhia de Jesus(安文思神父致澳门及广州诸耶稣会神父函)，ARSI, Jap.Sin(罗马耶稣会档案馆，和汉档)，162,fols.269-273v.}
\BookNote{72.}{〈耶稣会士冯秉正神父致同一耶稣会的法国传教区会长赫苍璧神父的信〉，页66。}
\end{BookNotes}

% 原书第 289 页（PDF 第 302 页）
\BookPage{289}{302}
\begin{BookQuote}
之子宋荦，年十四，入备宿卫，改官黄州通判，游至江苏巡抚是也。今满汉大员子弟，皆须十八岁始准入仕，无童穉拜官之例。惟皇子及诸王侍从小臣中有曰「哈哈珠子」者，清语为幼男之称，名虽幼男，亦非年届十八岁弗用，仅存
\end{BookQuote}

其义而已。73

福格只追溯到康熙中期，说是当年未满十八岁也可以入仕，后来连皇子、诸王的哈哈珠子也不用小男孩，全都用十八岁以上的男子，以致于哈哈珠子名实不符。事实上，哈哈珠子的制度应该是清初以来，甚至入关之前就有的作法，以男童作为皇子或年轻皇帝的随从兼伴侣吧。《雍正上谕内阁》收录了康熙六十一年十一月十七日的一件上谕，其文云：

\begin{BookQuote}
下五旗诸王属下人内，京官自侍郎以上，外官自州牧县令以上，该王辄将其子弟挑为包衣佐领下官及哈哈珠子、执事人折挫使令者甚众。嗣后着停止挑选。其现在行走人内，系伊父兄未任以前挑选者，令其照常行走。若系伊父兄既任以后挑选者，俱着查明撤回。或有过犯，该王特欲挑选之人，着该王将情由奏明，再行挑选。74
\end{BookQuote}

这段资料说明了下五旗王公也有将属下人员子弟挑选为哈哈珠子或执事人（拜唐阿、柏唐阿、栢唐阿）的惯例。虽然不是针对上三旗而发，那是因为上三旗的旗主正是皇帝本人，而皇帝当然也有挑选上三旗属下担任哈哈珠子等的作法。本条上谕也意味着哈哈珠子的职务是相当低下的工作。75

康熙皇帝本人出生于顺治十一年（1654），顺治十八年（1661）即位时虚岁八岁。他的哈哈珠子与他的年纪，所差应不至太大。

赵昌的生年，依据耶稣会士冯秉正在1731年10月10日寄给法国赫苍璧神父的

\vfill
\begin{BookNotes}
\BookNote{73.}{福格，《听雨丛谈》（北京:中华书局，1959、1984)，页251-252。}
\BookNote{74.}{收在中国第一历史档案馆编，《雍正朝汉文谕旨汇编》（桂林市：广西师范大学出版社，1999)，第6册，页10。康熙皇帝在四天前才驾崩，不知雍正皇帝为何急着发表这样的一道上谕！}
\BookNote{75.}{《清代内阁大库散佚满文档案选编》（天津：天津古籍出版社，1992)，页6-7收录的文件涉及康熙二年，一位年仅十二岁的哈哈珠子常保因为在宫庭院内打扫完毕之后，在廊下睡觉被发现，因而被判处鞭责一百下的惩罚。可见得哈哈珠子的性质。参考唐英凯，犹龙柱下窥藏史，呼凤舆前听诵诗一读《清代内阁大库散佚满文档案选编》第3、4册〉，《清史研究》，1994年第1期，页96。}
\end{BookNotes}

% 原书第 290 页（PDF 第 303 页）
\BookPage{290}{303}
\noindent 那封信函的内容推算，雍正七年时赵昌73岁76，则他大概出生于1657年，比皇帝小了三岁。此外，康熙朝的名人曹寅出生于1658年，比康熙皇帝小四岁；李煦出生于1655年，比皇帝小一岁。曹寅与李煦皆可能是康熙皇帝的哈哈珠子。

除了这几位之外，另有一人也可能，他是民初上海有名的文人杨钟羲的先人。杨钟羲本人及家族在清朝覆亡以前为内务府正黄旗管领下人。杨钟羲在其《来室家乘》一书中，很详细地记载了他的家族谱系。他说：

\begin{BookQuote}
尼堪杨氏家世纪略。吾家先世居辽阳地方。天聪二年隶满洲都统内务府正黄旗头班管领。始祖讳讨塞（《八旗满洲氏族通谱》作陶色）从龙入关，官正黄旗内管领，三任织造。妣关氏，无所出，以杰臣公（官司库)、华塞公（《通谱》作华色，袭世职内管领，河东巡盐御史）为养子。世祖章皇帝以讨公之未育也，乃出宫人邰氏、苏氏为之配，实生搬臣公、摇臣公（搬臣公赐名搬不动，官膳房，蓝翎侍卫；摇臣公赐名摇不动）。华公无子，以摇臣公次子和善公袭内管领（《通谱》作浩善），掌关防管理三旗内管领事务。乾隆三十九年卒。…康熙二十四年八月十三日，和公生，袭内管领，掌关防，管理三旗内管领事务。…乾隆十八年九月二十九日卒，年九十。77
\end{BookQuote}

《来室家乘》关于和善（浩善）的卒年似乎有误。他果若生于康熙二十四年为实，即生于1685年，则1774年方为九十岁。1774年即乾隆三十九年。和善的生身之父为摇臣，从和善的生年往后倒推三十年（假定他三十岁时生和善），摇臣自己可能出生于1655年左右，年纪正与康熙生年相当！虽然是因为摇臣和他的兄长搬臣名字的关系，才被分别赐名为摇不动和搬不动，不过，感觉上这很像是儿童相互戏要时取的绰号，而不像是庙堂之上的正式赐名。据此推断，搬臣和摇臣也可能是康熙幼年或青少年时的哈哈珠子。

\vfill
\begin{BookNotes}
\BookNote{76.}{〈耶稣会士冯秉正神父致同一耶稣会的法国传教区会长赫苍璧神父的信〉，页70。}
\BookNote{77.}{《来室家乘》原发表于《中和月刊》。杨家骆将之改题为〈雪桥自订年谱〉，收入《雪桥诗话》(台北:鼎文书局，1971)，第1册，此见页1-3。参考郁辉，〈论满族家谱《来氏家乘》的纪年方式〉，《满族研究》，2009年第1期，页64-67。}
\end{BookNotes}

% 原书第 291 页（PDF 第 304 页）
\BookPage{291}{304}
康熙皇帝即位时虚岁八岁。他在康熙六年（1667）亲政。康熙八年他主导了一场「擒鳌拜」的大戏。有趣的事，康熙擒鳌拜既然是一件重大的历史事件，应该会有不少研究吧。实际上却不多。主要的原因是文献不足之故，而挽近最主要的发现应该是推论出鳌拜本人被捕的时间是康熙八年五月初十日。78

擒鳌拜一事，在康熙朝之后甚久才有人笔之于书，一则是昭（1776-1829）的记述，提到以「羽林士卒」擒之。79生年与昭楗相同的安徽桐城人姚元之（1776-1852）则在其所着《竹叶亭杂记》一书中记载：

\begin{BookQuote}
圣祖仁皇帝之登极也，甫八龄。其时大臣鳌拜当国，势焰甚张，且以帝幼，肆行无忌。帝在内，日选小内监强有力者，令之习布库以为戏（布库，国语也，相斗赌力）。鳌拜或入奏事，不之避也。拜更以帝弱且好弄，心益坦然。一日入内，帝令布库擒之，十数小儿立执鳌拜，遂伏诛。80
\end{BookQuote}

昭与姚元之的生年距离康熙擒鳌拜的时间已经相差一百多年，所记录的内容不免属于辗转耳食之听闻，时移事迁，细节之处更不可能准确。「羽林力士」与「内监强有力者」的说法之外，「哈哈珠子」其实也很可能才是这批受命动手的人，他们是少年男子，不是内监。果若如此，则这些「哈哈珠子」当中就可能包括前面提到的曹寅81、李煦。擒鳌拜时，康熙十六岁、曹寅十二岁、李煦十五岁。还有杨钟羲的先人搬臣与摇臣，康熙皇帝称为搬不动及摇不动，故而有可能是与皇帝练习「布库」，皇帝随口叫出来的绰号。至于赵昌，当时他年纪为十三岁，也非常可能是其中的一员。

\BookSubheading{四、赵昌的各种职称}

即使赵昌担任哈哈珠子侍童的角色没有疑问，以他出生于1657年来推算，他要到1674年才十八岁大，也就是康熙十三年以后他才可能被当成大人来看待。陈青

\vfill
\begin{BookNotes}
\BookNote{78.}{白新良，〈康熙擒鳌拜时间考〉，《满族研究》，2005年第3期，页74-77。}
\BookNote{79.}{昭，《啸亭杂录》（北京：中华书局，1997），页5。}
\BookNote{80.}{姚元之，《竹叶亭杂记》(北京:中华书局，1997)，页1。}
\BookNote{81.}{参考刘上生，〈曹寅入侍康熙年代考〉，《中国文学研究》，2000年第2期，页26-29。}
\end{BookNotes}

% 原书第 292 页（PDF 第 305 页）
\BookPage{292}{305}
\noindent 松的文章引用顾卫民的《中国天主教编年史》，说康熙八年十一月十六日，皇帝派遣礼部大员佟国刚和侍卫赵昌前往汤若望墓前致祭。顾卫民所据为《熙朝定案》。但是在韩琦所编订的《熙朝崇正集·熙朝定案（外三种）》相关的部份，只提到礼部官员，完全没有提到赵昌。陈青松推测这应如韩琦所言「《熙朝定案》事实上是一部没有定本的著作」，也就是说顾卫民有不同的版本可用。然而韩琦除了依据徐家汇藏书楼、巴黎法国国家图书馆、罗马耶稣会档案馆等所藏版本校过之外，还参考过其他多种异本，应该极为齐全了。82因此，陈青松提到赵昌在康熙八年被派去汤若望墓前致祭的说法，就有待进一步的证据来支持。再者，如前所言，康熙八年时，赵昌虚岁只有十三岁，可能才开始在康熙身边当差不久，他的年纪、身分也与佟国刚不相称，因此顾卫民和陈青松所言，或许应稍作保留。

不过，至迟到康熙二十一年时，赵昌的名字已经正式出现在中文史料中。这里指的是一道上谕，而该上谕刻在传教士利类思(LudovicoBuglio)。的墓碑上。该上谕的文字如下：

\begin{BookQuote}
谕南怀仁（FerdinandVerbiest）等：今闻赵昌来奏，利类思年老久病，甚是危笃。朕念利类思自世祖章皇帝时至于如今，效力多年，老诚质朴，素知文翰。□□等俱系海外之人，利类思卧病京邸，绝无亲朋资助，深为可悯，□特赐银二百两、缎十疋，以示朕优念远臣之意。特谕。康熙二十一年九月初七日。83
\end{BookQuote}

19世纪末黄伯禄所编辑的《正教奉褒》一书也收录了该文件，并且接着说：「初八日，上遣一等侍卫噶，同侍卫赵昌、龚、萨来堂，赐茶酒祭奠。」84从康熙二十一年奉派奠祭利类思一事来看，赵昌为皇帝服务而与西洋人往来，一定已经有

\vfill
\begin{BookNotes}
\BookNote{82.}{陈青松，〈赵昌家世及其与传教士的往来—兼述其在康雍时期的际遇〉，页8及注16；顾卫民，《中国天主教编年史》（上海：上海书店出版社，2003)，页184；韩琦、吴旻校注，《熙朝崇正集·熙朝定案（外三种)》，页4,12-13,83.}
\BookNote{83.}{高智瑜、马爱德(编)，《虽逝犹存—栅栏：北京最古老的天主教墓地》（澳门：澳门特别行政区政府文化局，2001)，页144-145。}
\BookNote{84.}{《熙朝崇正集·熙朝定案(外三种）》，页338-339。黄伯禄的书未注明出处。韩琦发现他大量参考了《熙朝定案》。韩琦也指出《熙朝定案》错将二十一年误作二十二年，我们发现《正教奉褒》已经改正过了。}
\end{BookNotes}

% 原书第 293 页（PDF 第 306 页）
\BookPage{293}{306}
\noindent 一段时间了。有趣的是赵昌以外的其他三名侍卫（噶、龚、萨）都没有被记录较为完整的名字，以致于难以进一步查考。其后赵昌扈从南巡，经常以代表皇帝的名义造访地方上的天主堂，赏赐物品。85

虽然赵昌在中文史料中明确留下纪录已经是康熙二十一年的事，有机会贴身观察宫中活动的西洋传教士倒是早早就记录了这个人的若干事迹，大都也是一位侍从人员所做的各小事，详情可参考金国平的文章。另一方面，金国平以中文史料为中心，对他历年的职务与职称作了相当完整的整理，以下就以他的作品为基础，稍行补充，简述一下赵昌在康熙朝的种种职称（包括某些颇为模糊的称谓）和他所从事的一般工作。

金国平整理了文献上所见赵昌的职称包括:康熙二十一年，「御前」；康熙二十五年，「侍卫」；康熙三十三年，「养心殿监造六品官兼景山火器总管」；康熙四十一年，「御书处赵昌」86；康熙四十五年，「养心殿总监造」；康熙四十七年:「养心殿看制造等项赵昌」87；康熙四十八年，「养心殿（赵昌）」、「养心殿总管」；康熙五十一年、五十二年，「养心殿监造」；康熙五十三年，「御前大臣」、「养心殿校对六品官」；康熙五十七年:「景山内御制威远将军8⁸总管，景山炮、鸟枪监造」、「养心殿监造六品官兼景山火器总管」。以上除康熙五十三年的「御前大臣」外，金国平皆有清楚出处。不过，因为部份史源为《康熙朝满文朱批奏折全译》，而该书满译汉的用字没有统一，因此未必代表当康熙时期的正式汉名职衔，不过也相去不远就是了。89

康熙五十三年有文献称赵昌为「御前大臣」之处，金国平并未提供出处，无法

\vfill
\begin{BookNotes}
\BookNote{85.}{这在《正教奉褒》中搜罗了很多的记录，而孟德卫（DavidE.Mungello)的书中也常加引用。参考 David E. Mungello, The Spirit and the Flesh in Shandong, 1650-1785, pp. 68-71 ; David E. Mungello, The Forgotten Christians of Hangzhou, Pp. 60-63.}
\BookNote{86.}{Antonio Sisto Rosso, Apostolic Legations to China in the Eighteenth Century, p. 234 所收录原藏梵蒂冈教廷图书馆文件。}
\BookNote{87.}{Antonio Sisto Rosso, Apostolic Legations to China in the Eighteenth Century,p. 250 所收录原藏法国国家图书馆文件。}
\BookNote{88.}{按、「威远将军」为特定等级的大砲名称。}
\BookNote{89.}{参考白新良，〈已刊康熙朝满汉文奏折正误〉，《清史研究》，2003年第4期，页74-80。}
\end{BookNotes}

% 原书第 294 页（PDF 第 307 页）
\BookPage{294}{307}
\noindent 查回。恐怕是一种不精确的纪录、或是客气的称谓。同一年有一件满文奏折提到赵昌，他的职称译作「养心殿校对六品官赵昌」。90因为赵昌是皇帝身边的人，常被称为「大人」。不过，六品官理应不可能有「御前大臣」的正式称呼。

至于康熙二十一年的「御前」，金国平指出最早见于《熙朝定案》。91他并且推论说「御前」就是「御前侍卫」的简称，并且征引他人著作说「御前侍卫」是二品武官。这个论说不正确。事实上，《熙朝定案》等文献所称之「御前」，其实都只能说他在近处侍御，在皇帝左右任职而已。若说他是侍卫，大概也顶多像他的祖父一样，是个「蓝翎」吧。92

除此之外，金国平引用康熙五十五年以殊色印刷送往欧洲的「红票」内文，将「武英殿等处监修书官赵昌」当成是赵昌的职衔之一。此一写法颇令人误会赵昌也担任武英殿官员。事实上「红票」前头的衔名为「武英殿等处监修书官伊都立、王道化、赵昌等」。所谓「武英殿等处」指武英殿和养心殿，伊都立为武英殿总监造（也称监修书官）93，但是赵昌和王道化两人却都是养心殿（造办处）的主管。

收在《康熙朝满文朱批奏折全译》，康熙四十五年正月二十二日，〈广储司银皮库奏报装饰养心殿所用物件数目折〉提到相关主管官员就有「武英殿总监造员外郎赫世亨」、「养心殿总监造赵昌」、「养心殿总监造王道化」等人。94员外郎为从五品官，而各项资料显示赵昌为六品官，因此武英殿总监造赫世亨与赵昌一起出现在文件上时，赫世亨总是列名在前。

\begin{BookQuote}
《康熙朝满文朱批奏折全译》不止一处，提到赵昌都是「养心殿总监造赵
\end{BookQuote}

\vfill
\begin{BookNotes}
\BookNote{90.}{康熙五十三年十一月二十五日，〈内务府奏请将委校对刘文林等并入派差内分份银折〉，《康熙朝满文殊批奏折全译》，页985-986。}
\BookNote{91.}{按、他处（如在《正教奉褒》一书）也多次见到。}
\BookNote{92.}{昭，《啸亭续录》，卷一，〈花翎蓝翎定制〉条云：「凡领侍卫府官、护军营、前锋营、火器营、銮仪卫满员五品以上者，皆冠戴孔雀花翎，六品以下者冠戴鹖羽蓝翎，以为辨别。」，见前引《啸亭杂录》，页382；福格，《听雨丛谈》，卷一，〈花翎〉条也有相关叙述，见页9-10。参考秦国经，〈清代宫廷的警卫制度〉，《故宫博物院院刊》，1990年第4期，页64-71。}
\BookNote{93.}{「武英殿监修书官伊都立」亦见于《明清时期澳门问题档案文献汇编》，页113。他在赵昌落难之初，与王道化一起列名在查抄赵昌家产的内务府官员名单中。参考《雍正朝满文朱批奏折全译》，上册，页2。}
\BookNote{94.}{《康熙朝满文朱批奏折全译》，页405。}
\end{BookNotes}

% 原书第 295 页（PDF 第 308 页）
\BookPage{295}{308}
\noindent 昌」或其替代性的称法「监造」、「监督」之类。此外，金国平也提到康熙五十年（1711），《佩文韵府》书成，赵昌以「养心殿总监造」的职衔列名于纂修、监造官员职名之中。

此外，赵昌的名字也出现在康熙五十年代的其他官修书籍。康熙五十二年(1713)，举行六旬万寿盛典。同年开始编辑《万寿盛典初集》一书，赵昌的名字也在书中出现。

康熙五十五年（1716），《万寿圣典初集》完成付梓，奉旨开列有功人员职名。收在安徽美术出版社复印的《万寿盛典初集》为武英殿当时的原刊本”，赵昌以「养心殿总监造六品官」名义与「养心殿总监造王道化」两人同以「监造官」的职衔列名于全部人员之末尾。96乾隆年间编修的「四库全书」也收录了《万寿圣典初集》一书。不过，奇怪地是在其书首，「纂修职名」下，康熙五十五年四月二十八日奉旨开列的人名最后一项一样是「监造官」，然而列出的名字却比武英殿版多出四位，即：

监造官：

\begin{BookQuote}
养心殿总监造六品官臣赵昌、养心殿总监造臣王道化、武英殿总监造兼会计司员外郎臣伊都立、武英殿总监造兼佐领臣张常住、武英殿总监造臣李国屏、原任云南巡抚臣佟毓秀。97
\end{BookQuote}

佟毓秀也是个画家，他在康熙四十四年(1705）时任云南巡抚8，在编辑《万寿圣典初集》之前，因此在《万寿圣典初集》编成时当然可以称作「原任云南巡抚」。其他伊都立、张常住与李国屏三人在康熙五十五年时都是内务府的现任官。因此，「四库全书」本的《万寿圣典初集》有可能是依据原始的稿本抄写的，而非依据武英殿版的刻本。令人好奇的是武英殿版为何要略去这四个人的衔名呢？还

\vfill
\begin{BookNotes}
\BookNote{95.}{按、安徽出版社所收《万寿盛典初集》仅选录数卷，并不完整。「四库全书」本之《万寿盛典初集》才是完整版。}
\BookNote{96.}{王原祁奉敕纂修，《万寿盛典初集》（合肥：安徽美术出版社，2002)，页（29)65-66。}
\BookNote{97.}{王原祁、王奕清等奉敕撰，《万寿盛典初集》（台北:台湾商务印书馆，「文渊阁本四库全书」），页（653）18。}
\BookNote{98.}{参考他的奏折，见《康熙朝汉文殊批奏折汇编》，页237。}
\end{BookNotes}

% 原书第 296 页（PDF 第 309 页）
\BookPage{296}{309}
\noindent 有，六品官的赵昌为何在全部「监造官」中位居首位呢？这个问题或许有趣，但为免离题过远，暂且搁下，以俟他日。

事实上，伊都立、张常住与李国屏三人任职武英殿，恐怕才是真正的「监造者。赵昌任职于养心殿，但因常侍皇帝，经常奉命传旨。据《万寿圣典初集》，卷四十，康熙五十二年王原祁的奏折所云，当年「五月初三日，养心殿监造赵昌等传旨，发下兵部右侍郎宋骏业所画万寿图稿，命臣绘画。」即是他所担任的真正角色。99

根据以上资料，从康熙三十三年前后以迄六十一年，赵昌的职务都是养心殿的监造官，100品级为六品。他若曾经是蓝翎侍卫，该职一般也是六品。赵昌大半辈子的官品都是六品。101果尔，这便是一个有趣的问题了！如资料所示，他至少在康熙六十一年以前倒算回去的三十余年间，乃至于更久，他的官品都没有调升过！

赵昌的主要工作，或说真正的职称是养心殿造办处的监造（亦称「总监造」）。养心殿造办处的工作在于「成做活计」，也就是制造宫中的用物与工艺品。很可能在康熙八年以后，养心殿就开始「成做活计」。不过，要到康熙三十年才正式成立养心殿造办处，而在三十二年以后才正式成立各作（各个作坊）。102这也是为什么要到康熙三十三年左右才看到赵昌带着「养心殿（总）监造」的职衔之故。赵昌也的确实际负责监督和管理这方面的工作。103

不过，赵昌近侍皇帝，因此也经常口啣天宪，代传旨意。此外，他也时而被派予兼任其他的任务，包括随皇帝南巡，代替皇帝到教堂行礼之类。有很长的一段时间，他可能与宫中的火器管理有关，陈青松、金国平的文章都论及此点。雍正皇帝

\vfill
\begin{BookNotes}
\BookNote{99.}{见《万寿盛典初集》（安徽美术出版社)，页(29)582-583。亦见王原祁、王奕清等奉敕撰，《万寿盛典初集》（台湾商务印书馆)，页（653)460。}
\BookNote{100.}{「养心殿监造赵昌」一称也见于《明清时期澳门问题档案文献汇编》，页100及同书多处。}
\BookNote{101.}{参考，陈青松，〈赵昌家世及其与传教士的往来—兼述其在康雍时期的际遇〉，页6。}
\BookNote{102.}{参考郭福祥，〈康熙时期的养心殿〉，《故宫博物院院刊》，2003年第4期，页30-34；吴焯，〈来华耶稣会士与清廷内务府造办处〉，《九州学林》二卷二期（2004年，夏季)，页65-101。}
\BookNote{103.}{相关事实可参考《康熙朝满文殊批奏折全译》，页225-249，〈内务府奏报宫内各处饭房裁减腌鲜各菜数目折〉；页317，〈广储司银库奏报装饰所用物件数目折〉；页985-986，〈内务府奏请将委校对刘文林等并入派差内分份银折等奏折。}
\end{BookNotes}

% 原书第 297 页（PDF 第 310 页）
\BookPage{297}{310}
\noindent 将他下狱后一段时间也提到他擅长火器。

雍正二年五月十二日，上谕云：

\begin{BookQuote}
…如放砲、放鸟鎗及扛抬鹿角者，俱系八旗汉军。从前内府佐领下那尔泰、赵昌善于放砲、放鸟鎗，皇考因交与那尔泰、赵昌管辖。乃伊等遂具奏将步兵砲手鄂尔博之缺，挑选内府佐领下人。近朕以内府佐领下人如何可占各旗分例，敕谕八旗：汉军步兵砲手鄂尔达之缺一概驳回各旗。...104
\end{BookQuote}

虽然有人以为这是赵昌获罪的理由之一，其实并没有拿来增加对赵昌的处罚，看来比较像是借题发挥。雍正皇帝把该事件提高到内务府包衣占用汉军八旗缺分的问题上，要求将砲手的缺分还给八旗汉军。无论如何，毕竟那是雍正皇帝的「皇考」命令赵昌去做的。限于篇幅，不宜多谈缺分的问题，但此一条上谕可说明赵昌对枪、𪿫的施放的确颇为娴熟。另一位被雍正皇帝点名的那尔泰，在康熙五十至五十五年前后，担任苏州织造衙门的乌林达（库掌）。105

赵昌的任职表现应该不错，很少看到他被皇帝斥责或处罚的事。金国平的文章引用了耶稣会士傅圣泽（Jean-FrançoisFoucquet)1716年的一封信件，该信件收藏于罗马耶稣会档案馆，归在「和汉档」之内。傅圣泽在该信中:「因此，出于利益，赵昌成为法国人的敌人。20多年前，皇帝下令把他痛打一顿，因为皇帝意识到他想加害他们。」由于欠缺进一步解释，关于此事的来龙去脉，无法获知其详。1716年之前二十年为1696年，即康熙三十五年。该年之前，赵昌已经在养心殿造办处任职，他或许因业务过失受责。至于他是否是法国人的敌人，或许只是傅圣泽一己的想法而已。下一节将说明，在法国籍耶稣会士于1687年抵华之初，还是靠赵昌帮忙才得以留下来活动，而且后来赵昌也帮助法国传教士取得蚕池口的土地建立住院与教堂。傅圣泽的说法或许只能暂作参考。

既有研究发现康熙五十五年十一月，赵昌和养心殿的其他主管被废太子胤初之子弘皙要求，让养心殿工匠华色为他制造逾制的珐瑯火链（火镰，打火器）。事

\vfill
\begin{BookNotes}
\BookNote{104.}{中国第一历史档案馆编，《雍正朝汉文谕旨汇编》，第6册（上谕内阁），页129。}
\BookNote{105.}{见故宫博物院明清档案部编，《李煦奏折》（北京:中华书局，1976)，页91,191-193；乌林达那尔泰在康熙朝曾经单独进折或领衔进折一件。见参考白新良，〈康熙朝进折人员考略〉，页40。}
\end{BookNotes}

% 原书第 298 页（PDF 第 311 页）
\BookPage{298}{311}
\noindent 发，工匠被流放，而赵昌也与养心殿的另一位监造王道化一并受到处罚。106这一件事与他的养心殿总监造职务有关，但情非得已。

除此之外，就现有资料看来，作为皇帝身边小臣的赵昌在康熙朝并没有犯过重大的过失。有的话，应该也只是职务上的小缺失吧，因此也就遭到过一些小处分，但是应该都不严重，所以在康熙皇帝六十大寿庆典期间，他也获得「开复先前处分」的恩遇。康熙五十二年三月二十八日，吏部奉旨开列获此「恩典」的各项人员名单，排在第一位的是内务府总管赫奕，第二位是江宁巡抚张伯行、第三位是四川巡抚年羹尧，而「养心殿监造赵昌」则列名倒数第二。107

赵昌的本职是内务府养心殿造办处的总监造。不过，他最重要的差使或任务，却是负责与西洋人往来。这当然是因为他常伴君侧，受皇帝倚重。有些学者说，他也兼有监视这些西洋人的任务。如金国平引1686年徐日升与另一神父信函所称，赵昌是皇帝的「眼睛、耳朵和嘴巴」108，换言之，皇帝的耳目与代言人。

\BookSubheading{五、赵昌与西洋人}

赵昌与西洋人事务发生关系，开始得很早。前已面经提过，康熙二十一年九月初七日（1682年10月7日），在给南怀仁的上谕中就提到「今闻赵昌来奏」的字样。当年赵昌应为二十六岁，显然早已受命为皇帝处理服务宫廷的西洋传教士的事情。

在北京或中国其他地方所能接触到的西洋人当然未必是传教士，但是总以天主教传教士为主。不过，有趣的是他也接待过荷兰使节团。在郑克坏降清、清廷开放海禁之后，荷兰人派遣了一个使节团前往北京。在那一次的接触中，赵昌的身影不断出现。卫思韩（JohnE.Wills,Jr.）说，在1686年（康熙二十五年）：

\begin{BookQuote}
从7月31日到9月17日间，荷兰使节一行留在北京。皇帝命令耶稣会士们与
\end{BookQuote}

\vfill
\begin{BookNotes}
\BookNote{106.}{杨启樵，〈刘心武先生的《揭秘红楼梦》质疑〉，页304。}
\BookNote{107.}{《万寿盛典初集》(台湾商务印书馆)，页(653)286-289。}
\BookNote{108.}{原注:Carta de 20.9.1686do P. Tomás Pereira aoP. Sebastião deMagalhães(徐日升神父致塞巴斯蒂昂·德.马加良埃斯神父函)。}
\end{BookNotes}

% 原书第 299 页（PDF 第 312 页）
\BookPage{299}{312}
\begin{BookQuote}
荷兰人保持密切的接触，而他们也就经常拜访荷兰人。他们有时带着赵昌一道去，这人是宫廷侍卫的一员，也是皇帝与耶稣会士之间的关键性说项者。赵昌偶尔也独自去拜访荷兰人，从而成为他们喜爱的一位居间人，借以打探宫廷对他们要求的决定动向，寻求有利的决定。他有时候用拉丁语与同吧芝交谈，据说对西方科学感兴趣。109
\end{BookQuote}

赵昌在这个1686年荷兰贡使事件中也从使节团方面收受了几座钟表当作礼物。110这点让我们想到他在雍正元年被抄家时，抄出来的东西当中，不就有「大小钟表六」件吗！

其后不久，法国国王路易十四（LouisXIV）派遣五位传教士以所谓的「国王的数家」的名义前往中国。为了避免葡萄牙人的阻挠，接近中国以后，改搭广东人王华士的商船，沿岸北行，抵达宁波。据说，耶稣会士洪若（JeandeFontaney）透过人在浙江的义大利籍传教士殷铎泽写信告知服务于北京宫廷的南怀仁他们抵达中国的讯息，而南怀仁则请赵昌协助，从而使得这几位传教士得以在抵达中国的第二年(康熙二十七年，1788）获准到北京觐见皇帝。111

\vfill
\begin{BookNotes}
\BookNote{109.}{John E. Wills, Jr., "Some Dutch Sources on the Jesuit China Mission, 1662-1687," Archivum Historicum Societatis Iesu,54(1985),pp.267-294.引文见p. 277。该文的中译见卫思韩,关于1662-1687年耶稣会中国传教团的一些荷兰史料〉，收在中国人民大学清史研究所编，《清史研究集》第7辑(北京:光明日报出版社，1990)，页352-377，特别是页360。此处的引文已经本文作者重译。卫思韩其实在更早的作品中就已经提到赵昌在1686年这次荷兰使节团事件中扮演的角色。详见 John E. Wills, Jr., Pepper, Guns, and Parleys: the Dutch East India Company and China, 1662-1681 (Cambridge, Mass.: Harvard University, 1984) , Pp. 204-212: John E. Wills, Jr., Embassies andIllusions: Dutch and PortugueseEnvoys toK'ang-hsi, 1666-1687(Cambridge, Mass.: Harvard University, 1984) , Pp. 160-165 and171。有关本次荷兰使节一行的完整研究见J. Vixseboxse, Een Hollandsch Gezantschap naar China inde Zeventiende Eeuw(1685-1687)(Leiden: E.J. Brill,1946)。引文中的吧芝即Vincent Paets，荷兰使团的正使，副使为LouisdeKeijser。卫思韩所据为荷兰东印度公司档案，在档案中赵昌被写作“Tsiaulauja(Chao-lao-yeh)”或接近的拼法，卫思韩经过与史景迁（JonathanSpence)、胡天龙（FrancisA.Rouleau)通信后确定这个人就是赵昌。}
\BookNote{110.}{参考前注。}
\BookNote{111.}{Isabelle Landdry-Deron, "Les Mathématiciens envoyés en Chine par Louis XIV en 1685," Archive for History ofExactSciences,no.55（2001)，pp.423-463；李景屏，〈康熙与路易十四施政之比较—从法国派遣科学传教团来华科考谈起〉，《文史知识》，2010年第3期，页36-43；覃波，〈康熙皇帝与法国远臣〉，页31。参考柯兰霓(ClaudiavonCollani)着、李岩译，《耶稣会士白晋的生平与著作》；德·博西耶尔夫人(Mme. deBossierre)着、辛岩译，《耶稣会士张诚—路易十四派往中国的五位数学家之一》（郑州：大象出版社，2009）。}
\end{BookNotes}

% 原书第 300 页（PDF 第 313 页）
\BookPage{300}{313}
康熙三十年，浙江兰溪、临安等县发生了反对天主教的行动。殷铎泽致函浙江巡抚张鹏翮，希望他出面制止，结果事与愿违，反而导致浙江省全面禁教的措施。殷铎泽于是透过服务于北京的传教士们设法。最后由张诚（Jean-François Gerbillon）、徐日升求见康熙皇帝。居间联络、转递消息的人就是赵昌。而赵昌不只是担任联络的角色，事实上还协助这些传教士们准备必要的文件、影响皇帝的心意。最后在康熙三十一年二月初三日（1692年3月20日），由礼部尚书顾八代领衔具题，初五日（1692年3月22日）奉旨「依议」，以正式的行政作业程序完成对天主教信仰与行教的包容。112有一位近代的天主教历史学者胡天龙（FrancisA. Rouleau）因此说：这之后的十年间，即在康熙皇帝对多罗事件作出反应之前，可以说是中国天主教的黄金时代。113顾八代的题本与奉旨的批红14也就顺理成章地被称作「容教诏令(Edictof Tolerance)。

面对这样一个欣喜的结局，人在北京的耶稣会士苏霖（JoséSuarez）写了一件详细的报告。我们摘引几个小段落以见赵昌的角色：

\begin{BookQuote}
（1691年12月21日）皇宫中有一个大臣，叫赵老爷，若称其为中国教会的功勋元老也不为过，对于处理我们有关的案件他是行家里手(？老手），有好多年的经验，他代我们向皇上讲明我们的请求，皇上知道了赵老爷与我们之间的友谊与忠诚，愿意把这个案子全权交给他处理。…皇宫里的赵老爷，凭着同样的口才与情感，再一次觐见皇上，请求审阅徐日升所带来的文件。…1692年1月1日，皇上派对我们有好感的也是皇上宠爱的赵老爷到我们住院来，随身带来了皇上用非常精彩的满文所写的一封信，当众读给我们听。...115
\end{BookQuote}

\vfill
\begin{BookNotes}
\BookNote{112.}{Antonio Sisto Rosso, Apostolic Legations to China in the Eighteenth Century, Pp. 128-129 ;张先清，康熙三十一年容教诏令初探〉，《历史研究》，2006年第5期，页72-86；赵昌的角色见页77-78。参考吴伯娅，〈康熙帝事考两则〉，《清史论丛（2009年号）》，页187-196，第一节「康熙容教令考」。}
\BookNote{113.}{Francis A. Rouleau, "Maillard de Tournon, Papal Legate at the Court of Peking: The First Imperial Audience (31 December 1705) ,"pp. 268-269.}
\BookNote{114.}{其文本可参考黄伯禄编，《正教奉褒》，页357-358。}
\BookNote{115.}{莱布尼茨，《中国近事—为了照亮我们这个时代的历史》（郑州：大象出版社，2005)，页1-36为梅谦立所译〈苏霖关于1692「容教诏令」的报告〉。以上的引文出现在页25-27。}
\end{BookNotes}

% 原书第 301 页（PDF 第 314 页）
\BookPage{301}{314}
\noindent 苏霖报告里所说的「赵老爷」就是赵昌。

1692容教诏令之后，天主教在中国的教务长足发展，可是涉及祭祖与尊孔的所谓「礼仪问题」又浮出台面。康熙皇帝与耶稣会士都支持祭祖与尊孔，其他天主教士则不然。耶稣会士李明（Louis-DanielLeComte）有一部著作就因此受到巴黎外方传道会的批评。他于是致函在北京的耶稣会士，请他们为中国礼仪问题提供一个信实可靠的答案。于是耶稣会士闵明我(ClaudioFilippoGrimaldi)、徐日升、安多(AntoineThomas）、张诚联名起草一份文件，请内务府员外郎赫世亨译为满文，由主事张常住116奏上。皇帝看后说:「这所写甚好，合乎大道。敬天地、事君亲、敬师长，是天下通义。就是这样，并无可改之处。」这是康熙三十九年(1700)的事。117赵昌似乎没有在这个事件中插上一脚。而将该文件译作满文的赫世亨却在稍后的「多罗事件」中，与赵昌一起担负起皇帝与这位教廷使臣之间的居间人角色。赫世亨（1645?-1708），完颜氏，时任武英殿总监造。在稍后几年间，与西洋人有关的事件交涉中，赵昌并没有缺席，但是赫世亨的角色显然不亚于他，在两人会衔上奏时，赫世亨的名字经常都排在赵昌前面。118其间的道理除了前述赫世亨的官品略高之外，赫世亨通常留京办事，而赵昌则往往随着皇帝出巡，随时往来北京与驻跸所在有关吧。119

\vfill
\begin{BookNotes}
\BookNote{116.}{两人均任职于内务府武英殿，见前。}
\BookNote{117.}{罗丽达，〈一篇有关康熙朝耶稣会士礼仪之争的满文文献一兼及耶稣会士的宣言书《Brevis Relatio》〉，《历史档案》，1994年第1期，页94-97。闵明我、徐日升、安多、张诚所起草的文件内容如下（罗丽达自满文回译成汉文）:「西洋学者所说中国礼俗有祭孔子及供奉祖先之礼，然而不解其中缘故。他们说：皇帝恩惠无所不至，其圣明睿智各国无不闻知。因此，有此等礼俗，必有缘故。我们请求将此缘故详细写明寄给我们。所谓祭祀孔子者，乃敬其传世之言之谓也。既然接受他的教诲，岂能没有跪拜之礼？因此天下之人皆崇拜孔子，尊以为师，即此心也。并非为求聪明智慧，或为官职俸禄而敬之也。」参考AntonioSisto Rosso, Apostolic Legations to China in the Eighteenth Century, pp. 140-145 ; Joseph Dehergne, S. J., "L'exposé des Jésuites de Pékin sur le culte ancêtres présenté à l'empereur K'ang Hi en novembre 1700," Actes du IIe Colloque International de Sinology, Chantilly,1977（Paris,1980），Pp. 185-229；博西耶尔夫人着、辛岩译，《耶稣会士张诚—路易十四派往中国的五位数学家之一》，页95-97。}
\BookNote{118.}{参考《康熙朝满文朱批奏折全译》，页419-435，487，两人会衔的各奏折。同书所收由两人共同具名上奏的最后一个奏折见页497，康熙四十六年三月二十七日的折件。事实上，在现存资料中，赵昌仅以「养心殿监督兆昌」（此处写作「兆昌」而非「赵昌」）的名义单独进折或领衔进折一件。见白新良，〈康熙朝进折人员考略〉，页40。有关多罗事件时赵昌与赫世亨的角色，可参考冯明珠，〈坚持与容忍—档案中所见康熙皇帝对中梵关系生变的因应〉一文。}
\BookNote{119.}{参考《康熙朝满文殊批奏折全译》所收的相关折件。}
\end{BookNotes}

% 原书第 302 页（PDF 第 315 页）
\BookPage{302}{315}
不过，在此期间，赵昌对于法国传教士取得皇城内蚕池口，原苏克萨哈之旧邸为住院，并于其旁起造天主堂（北堂）一事，也有所贡献。北堂于1699年开始兴造，1703年落成。120

虽然会衔上奏时赵昌排名在赫世亨之后，但在向多罗宣谕时，赵昌有时候往往是实际的执行者，而赫世亨只担任陪同的角色。如四十五年五月二十五日（1706年7月5日）所发生的事。传谕的内容还包括一句「赵昌尔糊涂也！」，康熙皇帝责备赵昌对照顾多罗的医生之处置不够妥善。121

康熙四十六年（1707），皇帝探用了准许传教士领取「印票」，宣示永留中国则可继续居留，否则便驱逐出境的作法。赵昌当然也参与了这些工作的细节。他与多罗直郡王胤禔、武英殿总监造赫世亨、张常柱…等人一起负责给票的事，在第一历史档案馆所藏的「内务府满文行文档」中留下相关的史料。122同时，为了厘清有关「中国礼仪问题」清宫与教廷双方面的态度，康熙皇帝数次派遣传教士去欧洲，但在长期等待下音讯全无。于是在康熙五十五年（1716）印制了所谓的「红票」，经由海、陆两道送往欧洲。123「红票」的内容，就是以内务府武英殿及养心殿官员传谕的方式来书写，所以起头就说：「武英殿等处监修书官伊都立、王道

\vfill
\begin{BookNotes}
\BookNote{120.}{参考杜赫德编、郑德弟译，《耶稣会士中国书简集：中国回忆录》第2册（郑州：大象出版社，2001），页1-10，〈在华耶稣会传教士杜德美神父致本会洪若翰神父的信（1704年8月20日于北京）。以上这个文件有一个英文版，见TheTravels of SeveralLearnedMissioners of the Society of Jesus… (London, 1714) , pp. 198-214, "A Letter from E. Jartoux, Missioner of the Society of Jesus, in China, to F. deFontenay, of theSame Society."关于蚕池口天主堂（北堂）的历史，参考查时杰，〈从蚕池口到西什库一天主教北京北堂的历史〉，收在天主教辅仁大学历史系编辑，《中梵外交关系史国际学术研讨会论文集》，页119-144。}
\BookNote{121.}{《康熙朝满文殊批奏折全译》，页419-420，〈武英殿总监造赫世亨奏报西洋人情形折〉。}
\BookNote{122.}{例见〈总管内务府为核查发给西洋传教士印票事致兵部咨文〉、总管内务府为颁给西洋传教士王义仁等印票事致礼部咨文〉、〈总管内务府为知照颁给印票与西洋人等名单事致礼部咨文〉各件汉译「内务府满文行文档」，收在《明清时期澳门问题档案文献汇编》，页73-82。「印票」也称作「信票」或简称为「票」。有关「印票」的讨论非常多，基本上可以参考AntonioSistoRosso, Apostolic Legations to China in the Eighteenth Century、《熙朝定案》、《正教奉褒》…等书，此不具引。}
\BookNote{123.}{汤开建、刘清华，（康熙时期艾若瑟出使罗马始末考（1707-1720)〉，收在李向玉、李长森主编，《明清时期的中国与西班牙国际学术研讨会论文集》（澳门：澳门理工学院中西文化研究所，2009），页233-280；王庆成，〈康熙朝「礼仪」问题的一件新发现文献一武英殿修书官寄西洋人书并考释〉，《社会科学战线》，1996年第6期，页204-205。}
\end{BookNotes}

% 原书第 303 页（PDF 第 316 页）
\BookPage{303}{316}
\noindent 化、赵昌等字寄与自西洋来的众人」，赵昌也列名在内。124最后，康熙五十九、六十年(1720-1721）罗马教廷再次派遣宗主教嘉乐（PatriarchCarloAmbrogio Mezzabarba)来华，赵昌也与其他武英殿、养心殿官员伊都立、李国屏、李秉忠等一道担任接待和传话的工作。125

从以上的叙述来看，至迟从康熙二十一年起，一直到康熙六十一年底皇帝驾崩为止，作为皇帝贴身小臣的赵昌，与西洋人的事务关系密切是不容置疑的。当然，他或许会施些小惠，帮帮传教士忙，但要说影响到皇帝的心意与决断，应该不大。

1704年2月12日（康熙四十三年正月初八日），当时仍是皇太子的胤初命令一位名叫陆伯嘉（Jacques Brocard）的传教士把几件铁制品涂成蓝色，法国籍传教士杜德美（PierreJartoux）在场，而白晋也在工作现场担任翻译。传教士们由于怀疑其中的一根鞭子可能为佛教法器，而坚持不肯执行皇太子的命令。虽然胤初强调那是兵器而非法器，白晋还是极力争辩，终于激怒了胤初。最后事情闹到皇帝那边，而结局是皇帝说他本人处事，不愿勉强传教士们做不愿意的事。他们于是向他磕了九个头，又去胤初宫门前做了同样的事，事情才告一了结。这中间「一位对宽容基督教的敕令贡献良多的赵（Tchao）姓官员」曾经站出来帮忙，传教士们（特别是白晋）才愿意承认对皇太子的不敬，事情才得以收场。126

作为皇帝与西洋传教士之间的居间人，赵昌总是要设法让西洋人知道皇帝的性情。在那样的理解与谅解之下，他才能着力去完成任务。当然，其间不免杂含着他个人的理解与推断。1704年，在处理以上事件时，赵昌告诉白晋等人：「你们是外国人，皇帝的仁慈是你们惟一倚靠。皇帝保护你们，永许你们传教，因为这种宗教是合乎道德的，它只让人们做合理之事。无论在宫中还是在外省，皇帝什么好处、

\vfill
\begin{BookNotes}
\BookNote{124.}{「红票」的图像参见 Antonio Sisto Rosso, Apostolic Legations to China of the Eighteen Century, p. 308或严嘉乐着，丛林、李梅译，《中国来信（1716-1735)》，页223-226。该文件的原稿见陈垣编，《康熙与罗马使节关系文书影本》，页29-30。}
\BookNote{125.}{相关中文史料可参考陈垣编，《康熙与罗马使节关系文书影本》，页40-66。研究作品可参考 Antonio Sisto Rosso, Apostolic Legations to China of the Eighteen Century, pp. 202-211, 329-330}
\BookNote{126.}{见1704年8月20日〈在华耶稣会传教士杜德美神父致本会洪若翰神父的信〉，参见注120。洪若翰即洪若，见前。}
\end{BookNotes}

% 原书第 304 页（PDF 第 317 页）
\BookPage{304}{317}
\noindent 什么荣誉没有给过你们啊！…」127以一种软中带硬的方式，诱导白晋等人作出某种程度的让步。但是对于不熟悉宫廷行事风格，乃至于态度上过度不必要强硬的人士，赵昌其实也会给予适度的忠告。例如，1706年时，赵昌就曾告诉多罗:「皇帝宽容（柔软的）缎子，但会粉碎（坚硬的）钻石。对抗过多会使您受到严厉对待；如善于屈从，您将使君主心软。」128可惜没有打动多罗。

不过，赵昌并不是对所有的西洋人都是那么友善，或者说不是所有西洋传教士都对他心存感念。在天主教人士之中，他一向接触的是耶稣会传教士，也给予他们支持。但在1710年，遣使会的传教士德理格（TeodoricoPedrini，1671-1746）自马尼拉进入中国，藉其音乐方面的才能，被送往北京服务宫廷。129此后，他便致力于打击耶稣会士。1715年，教宗格列孟十一世（Clement XI）发布了所谓〈自该日起〉（exilladie）的宗座宪章，禁止中国天主教徒保持祭祖与尊孔等礼仪。130次年（1716），该文件被送达广州。正好因公出差到广州的内务府官员李秉忠因此有机会获悉内容，并且奏知康熙皇帝。131而在同年11月初（康熙五十五年九月下旬）的某一天，罗马传信部的方济会教士康和之（CarloHoratij）132以北京副主教的身分，到耶稣会的教堂布达教宗的谕旨。而负责处理「有关问题的几位官员当时正好在法国教士的住所。」也就是在北堂（蚕池口教堂），结果皇帝下令将康和之逮捕，送入监狱。此时正好是发布「红票」的时间，也就是康熙五十五年九月之后。依据「红票」上的署名，当时管理西洋人事务的正是「武英殿等处监修书官伊都立、王道化、赵昌」等人。133捷克籍的耶稣会士严嘉乐在报导了以上事件之后，

\vfill
\begin{BookNotes}
\BookNote{127.}{同上注。}
\BookNote{128.}{〈耶稣会中国副省长安多神父寄往欧洲的备忘录〉，收在郑德弟译，《耶稣会士中国书简集》，第6册（郑州：大象出版社，2005），页145。括弧内文字为原译者所加。}
\BookNote{129.}{参考马国贤着、李天纲译，《清廷十三年：马国贤在华回忆录》（上海:上海古籍出版社，2004)页29-30；汤开建，〈明清之际西洋音乐在中国内地传播考略〉，《故宫博物院院刊》，2003年第2期，页46-54；刘亚轩，〈清初来华传教士德理格与中国礼仪之争〉，《衡水学院学报》，第10卷第6期（2008年12月），页47,72-75。}
\BookNote{130.}{AntonioSisto Rosso, Apostolic Legations to China of the Eighteen Century, Pp. 191-193 and passim.}
\BookNote{131.}{严嘉乐着，丛林、李梅译，《中国来信（1716-1735）》，页33。}
\BookNote{132.}{康和之，也写成康和子。他是义大利人，出生于Castorano，有些人因此误以Castorano为他的姓氏。参考Antonio Sisto Rosso, Apostolic Legations to China of the Eighteen Century, Pp. 274-275。}
\BookNote{133.}{王道化和赵昌的职衔都是「养心殿总监造」，在「红票」中被省略。}
\end{BookNotes}

% 原书第 305 页（PDF 第 318 页）
\BookPage{305}{318}
\noindent 特别提到德理格「在去年秘密呈交给皇帝的备忘录中造谣中伤，肆意诽谤耶稣会士和一位与我们友好的大臣。」134严嘉乐没有指出这位被德理格中伤的大臣是谁，但康熙皇帝直接告诉我们德理格指控的有两位，其中一位是王道化，另一位即是赵昌。刘亚轩进一步参考方豪的著作，说康熙五十九年元旦（1720年2月8日），德理格没有入宫向康熙皇帝行三跪九叩的贺礼，结果被捕受惩。刘亚轩说「赵昌借此机会，把德理格狠狠整治了一番。」135这一部份的叙述，方豪在其著作中其实只是翻译法文《中国遣使会志》(Memoires de la Congregation dela Mission en Chine)中与德理格有关的文字。依据该段文字的内容，赵昌虽然在事件发生过程中一度现身，然而没有证据可以说是他在进行报复。136马国贤（MatteoRipa）以他自己的角度记录这次的事件说：「那个姓赵（Chao）的官员是我们的死敌，他不等看见德理格做出些什么，就越格执行了皇帝的指令，下令说我们的伙伴应该马上被当作罪大恶极者对待。」马国贤前往中国，曾经得助于德理格的安排。137虽然据第一位来华的遣使会会士毕天祥(LuigiAntonioAppiani)的说法，这两个人「颇不相投，作风亦不同。」138不过，他们都敌视耶稣会会士。赵昌是养心殿总监造之一，负责「带西洋人」139，又与耶稣会士友善，因此被非耶稣会士们视为眼中钉也就不稀奇了。所以，即使赵昌在次日也想办法给予德理格这方面的人忠告，马国贤也只往坏的方

\vfill
\begin{BookNotes}
\BookNote{134.}{严嘉乐着，丛林、李梅译，《中国来信（1716-1735）》，页34；参考柯蓝妮（柯兰霓）着、余三乐译，〈纪理安—维尔茨堡与中国的使者〉，收在任继愈主编，《国际汉学》，第11辑（郑州：大象出版社，2004)，页152-173，尤其是页164。参考刘亚宣，〈清初来华传教士德理格与中国礼仪之争〉，页73。关于康熙五十四年德理格中伤赵昌及王道化一事，见于陈垣编，《康熙与罗马使节关系文书影印本》，页77文件及 AntonioSisto Rosso,Apostolic Legations to China of the Eighteen Century, p. 369.}
\BookNote{135.}{刘亚宣，〈清初来华传教士德理格与中国礼仪之争〉，页74。}
\BookNote{136.}{方豪，《中国天主教史人物传》，页470-471。}
\BookNote{137.}{参考马国贤着、李天纲译，《清廷十三年：马国贤在华回忆录》，页29-30。}
\BookNote{138.}{方豪，《中国天主教史人物传》，页469。}
\BookNote{139.}{康熙五十九年（1720)教廷大使嘉乐来华时，康熙皇帝的一件中文谕旨将武英殿、养心殿的总监造们作如下称呼：「养心殿武英殿等处管制造、带西洋人事伊都理、张长住、王道化、赵昌」。前二者为武英殿总监造，后二者为养心殿总监造。武英殿总监造和养心殿总监造在前此一、二十年间都负责西洋人事务，但此一汉文文件难得地把这件差使明文写在他们的职衔里。见《掌故丛编》，第2辑(1928)，页15b-16a;Antonio Sisto Rosso, Apostolic Legations to China of the Eighteen Century, pp. 329-330.}
\end{BookNotes}

% 原书第 306 页（PDF 第 319 页）
\BookPage{306}{319}
\noindent 面去揣测。140

事实上，德理格、马国贤与山遥瞻三人来到中国时，也为多罗带来晋升为枢机主教的文件及礼服。多罗在康熙四十九年正月十七日（1710年2月15日）正式接受。虽然他人被软禁在澳门，却也将此事上书报告给康熙皇帝知道。同时，他也将这三位神父推荐给康熙皇帝。康熙后来接受了这三个人，但是对多罗的奏折形式深感不满，赵昌免不了奉命传旨责问。141作为传话人的赵昌，奉命传谕责备多罗，自然引起德理格、马国贤这些人不快，显然也是人之常情。总之，马国贤、德理格等人的确有不喜欢赵昌的理由。无怪乎在记录赵昌被捕的事件时，马国贤对赵昌作出如下的总结：

\begin{BookQuote}
陛下在他的一份谕旨中宣布，他之所以要如此地惩罚赵昌，是因为他太傲慢，还滥用职权迫害欧洲人。所有这些，我只能把它归结为全能的天主的判决。这就是著名的赵昌的下场，他是铎罗枢机主教和所有天主教的公敌。142
\end{BookQuote}

多罗，近代研究者常写作「铎罗」。

六赵昌的结局

本文一开始时，为了考订赵昌的旗籍，已经简述过雍正皇帝即位没几天，就将赵昌拿捕枷号。当时对中、外造成了某种程度的震惊，也见诸传教士的报导。至于在中国人之间，当然就免不了议论纷纷了。

雍正二年正月初八日，宗室赖实奏报康熙帝驾崩之时有人胡乱议论的奏折指控当时驻扎鸟兰古木（在鸟梁海）的「殿后都统」图喇及副都统鄂三等人的不法情事。指控的内容包括议论中言及雍正皇帝甫一即位，即「追问阿哥，执拿大臣，以致甚为混乱」、「诸事均舅舅公隆科多、公总督年羹尧承办，伊等如何说即如

\vfill
\begin{BookNotes}
\BookNote{140.}{马国贤着、李天纲译，《清廷十三年：马国贤在华回忆录》，页85-86。}
\BookNote{141.}{陈垣编，《康熙与罗马使节关系文书影印本》，页2-3,15-16；AntonioSistoRosso,ApostolicLegations toChina of the EighteenCentury， Pp. 284-285；参考万明，〈意大利传教士马国贤与中西文化交流〉，收在黄时鉴编，《东西交流论谭（第二集）》(上海：上海文艺出版社，2001)，页50，注2。}
\BookNote{142.}{马国贤着、李天纲译，《清廷十三年:马国贤在华回忆录》，页106。}
\end{BookNotes}

% 原书第 307 页（PDF 第 320 页）
\BookPage{307}{320}
\noindent 何办，如同谕旨。若此等亲戚，则均可富贵。若属原用之人，则均治罪，籍没家产。」以及「执拿赵昌，籍没家产。抛弃李英贵，又籍没家产流放，更有讨厌之语。」143这些话显然都是在议论雍正即位初期的行为，而将赵昌及李英贵免官逮捕、抄家籍没都是议论的内容。

李英贵也写作李锳贵，内务府正白旗包衣旗鼓人，长期在内务府当差，康熙五十七年时也曾署理内务府总管，后来担任过户部侍郎。144《清实录》，雍正元年正月甲申（初四日）：「谕内阁，仓场侍郎李英贵患病，仓场事务着顺天府府尹陈守创暂行办理。」145展开对李英贵的清算。次年，李英贵被发往打牲乌拉。146

赵昌与李英贵均为内务府人，都在雍正继位后一、两个月内被逮捕，远在西北军中的流言讲的也是这两个人。不过，雍正皇帝下手的对象很多，只是时间稍比这两个人晚些罢了。

作为外国来的旁观者马国贤对赵昌的下场也有所观察，但是他对赵昌早有偏见，因此措辞颇重。

\vfill
\begin{BookNotes}
\BookNote{143.}{《雍正朝满文朱批奏折全译》，上册，页603-604。该奏折的翻译颇不理想，词不达意，因此奏折本文及殊批的内容其实不易掌握。不过，该奏折的主旨（不含内容的文字）以及完整的殊批，都以编辑过的典雅汉文收入《（世宗宪皇帝）上谕内阁》，卷十五，内容如下（汉文《上谕内阁》中图喇写成图腊、赖实写成赖士）：「（雍正二年正月）初八日，军前效力宗室赖士参奏驻防都统图腊、副都统鄂三与群小怨望诽谤圣躬，捏造拘拏诸大臣、凌逼众阿哥，纵恣隆科多、年羹尧擅权。等语。具奏。奉上谕：朕览尔此奏，不觉大笑。尔乃不知朕居心行事之人。若朕身有过愆，我圣祖皇考误将此位付托于朕，若辈虽不如此讥评，而朕返躬自问，亦觉赧然。若朕之品行尚堪承荷大业，能不负我皇考付托之重，有裨益于国家黎民，无论无此荒唐谬妄之语，即使有之，其议论之人，天必鉴之，徒自作孽耳！于朕何与？况此等琐屑怀私之事，朕实不为。若朕果以恩怨系心，肯为挟私报复之人，我皇考至圣至神，亦必不将大位传与朕躬也。且朕从来并无汲汲以居心行事表白于众之念，亦并不以此置之怀抱。不但尔等不能知朕，举世亦无知朕之人。今此一年内，中外之人亦略知大概矣。尔远在军前，是以尚未悉知；若再稍有所闻，自能豁然醒悟。俟尔豁然醒悟时，尔自览今日之所奏，亦必发笑也。尔亦不必愧祚畏惧，但当安心黾勉力赎罪，不必营营，枉劳心虑也。亲笔书此，特行晓谕。尔身为宗室，故如此恺切陈奏，朕心嘉悦。尔但宽心，勿为朕虑。朕必不至大误极谬，可以自保。又奉上谕：将赖士此奏，令两营大臣官员俱行观览，并朕朱批谕旨亦令观之，使众各一笑。」见《雍正朝汉文谕旨汇编》，第6册（上谕内阁），页100。}
\BookNote{144.}{台湾商务印书馆影印文渊阁「四库全书」本《八旗满洲氏族通谱》，页（456）307-308；《康熙朝满文殊批奏折全译》，页1325，〈内务府奏请庄头子弟拓荒交粮折〉。}
\BookNote{145.}{《大清世宗宪（雍正）皇帝实录》(台北:华文书局，1964)，3/28a（页54)。}
\BookNote{146.}{《雍正朝满文殊批奏折全译》，上册，页664-667。}
\end{BookNotes}

% 原书第 308 页（PDF 第 321 页）
\BookPage{308}{321}
\begin{BookQuote}
康熙皇帝死后不几天，上述的葬礼还在进行的时候，继位皇帝雍正颁布了一条谕令，以一种让整个帝国感到吃惊的方式显示了他的掌权。在他的命令下，官员赵昌被抓了起来，戴上了重重的镣铐，宣判以枷刑处死。木枷是一种戴着走的板子，有200磅重。这位傲慢朝臣的财产被没收充公，家人受罚为奴，妻妾被分配给他人。147
\end{BookQuote}

关于赵昌被捕及用刑，捷克籍的传教士严嘉乐在其押上1723年10月14日的书信中也提到：

\begin{BookQuote}
「枷」这种刑具是两块连在一起的木板，当中有一个圆洞，套在犯人的脖子上。因为枷很重很宽，可怜的戴枷的犯人只能一直躺在地上。他的手够不着嘴巴，所以吃饭得让人喂。很少戴枷的犯人能活过三个月，可这位忠心耿耿的人却熬过了这一关—皇帝后来赦免了他，只没收了他的田产。148
\end{BookQuote}

1723年10月14日为雍正元年九月十六日，赵昌被捕到这个时候，已经超过十个月了。据严嘉乐所述，他戴枷的时间可能在三个月左右，而非更久的时间。不过，皇帝并没有赦免他。赵昌终究还是被关入牢狱，直到死亡。倒是马国贤报导说他原先戴的枷重达200磅，是否属实呢？事实上，颇有可能。

雍正二年五月间发生一个「李元玺案」，雍正皇帝深恶此人，行刑者揣摩上意，结果他被戴上了三百斤重的枷！雍正二年十一月初六日，

\begin{BookQuote}
谕刑部：从前因李元玺等俱系恶棍，朕因曾降谕旨，着用大枷枷号。乃阿尔松阿、马尔齐哈居心奸险，欲以非刑之名加朕，竟用三百斤之枷枷号。刑部从前用过之枷，俱有斤数，交与刑部查明，仍照从前刑部用过大枷斤数将李元玺等一并换枷枷号。149
\end{BookQuote}

看来李元玺戴着三百斤的重枷近半年之久！至于赵昌所戴之枷若为200磅，则约为三百斤之半，算起来也接近一百公斤，已经是严重的酷刑了。他能一戴戴了三

\vfill
\begin{BookNotes}
\BookNote{147.}{马国贤着、李天纲译，《清廷十三年：马国贤在华回忆录》，页106。}
\BookNote{148.}{严嘉乐着，丛林、李梅译，《中国来信（1716-1735)》，页41。}
\BookNote{149.}{《雍正朝汉文谕旨汇编》，第6册(上谕内阁)，页165。}
\end{BookNotes}

% 原书第 309 页（PDF 第 322 页）
\BookPage{309}{322}
\noindent 个月，实在不容易。

另有一点是雍正皇帝大概无意直接下令杀死赵昌，理由也是不想背负那样的名义。所以雍正二年六月二十一日，当刑部建议将宗室苏努一案150的要角，皇九子胤糖的岳父、曾任浒墅关监督的七十（「七十」是人名151）处死时，有上谕云：

\begin{BookQuote}
七十甚属妄乱，应即行正法示众。但朕自御极以来，除强盗、故杀、谋杀人犯外，并未降旨特杀一人。七十虽应处斩，亦不足以当朕之初命正法之人。…七十着从宽免死，佥妻子发往三姓地方。…152
\end{BookQuote}

类似的考量，或许可以解释为何赵昌在雍正即位之初即已被捕下狱，却还能撑到雍正九年才亡故的原因。

末了，我们还是得提一下赵昌最后受洗一事。冯秉正的说法是说赵昌因为有姬妾数房，不合天主教一夫一妻制度的规矩，因此始终不能领洗，并且说「我本来在十年前就应该接受这种圣宠了。」153十年之前就是他被捕下狱的康熙六十一年。换言之，最后取名「若瑟」（Joseph）入教的赵昌，动念接受天主教信仰的时间不能说极早。但是回顾他在康熙二十八年曾经表示「虽然他经常造访天主堂，他对他所见的一切一概不明了。」154赵昌最后的改信，也还真的是一个颇大的转变。

\vfill
\begin{BookNotes}
\BookNote{150.}{关于苏努一家信奉天主教及受到雍正皇帝迫害一事，研究成果颇多，可参考方豪，〈清代旗人之信奉天主教与遭禁〉页1-23；冯佐哲，〈清宗室苏努举家信奉天主教〉，《紫禁城》，1990年第1期，页8-9；吴伯娅，《康雍干三帝与西学东渐》，页183-188；吴伯娅，〈苏努研究〉，收在张先清编，《史料与视界一中文文献与中国基督教史研究》，页61-82。}
\BookNote{151.}{参考《康熙朝满文朱批奏折全译》，页1268-1272，康熙五十六年十一月二十二日，〈内务府总管允陶等奏销养心殿等工程银两折〉；《雍正朝汉文谕旨汇编》，第6册（上谕内阁），页34-36, 131-133。}
\BookNote{152.}{《雍正朝汉文谕旨汇编》，第6册（上谕内阁），页138。}
\BookNote{153.}{〈耶稣会士冯秉正神父致同一耶稣会的法国传教区会长赫苍璧神父的信〉，页72。}
\BookNote{154.}{参考注34。}
\end{BookNotes}

% 原书第 310 页（PDF 第 323 页）
\BookPage{310}{323}
